\chapter{拉格朗日子流形}
\label{ch:lagrangian-submanifolds}

本章引入\emph{拉格朗日子流形}这一概念，把它看作最小作用量原理的\emph{自然边界条件}。遵循物理学家的术语，我们一般把由最小作用量原理产生的方程称为给定泛函的\emph{运动方程}。

为引出以下讨论，我们从余切丛 $M=T^*N$ 上带有其典范辛形式的最小作用量原理讲起。

\section{余法丛}
\label{sec:conormal-bundles}
\stfhIndex{I185}{la^ge^lang^ri^zi^liu^xing}{拉格朗日子流形（Lagrangian submanifold）}
\stfhIndex{I209}{zi^ran^bian^jie^tiao^jian}{自然边界条件（natural boundary condition）}

在余切丛 $T^*N$ 上，哈密顿作用量泛函定义为积分
\[
  \mathcal A_H(\gamma)=\int_\gamma\theta-\int_0^1H(t,\gamma(t))\,dt.
\]
通常把 $\theta$ 写作 $p\,dq$；我们也会不时采用这一写法。下面推导 $\mathcal A_H$ 的第一变分公式。沿给定道路 $\gamma$ 取变分 $\xi$，其支集包含在 $[0,1]$ 的某个小开子集中，并且相应的像包含在 $T^*U$ 中，其中 $U\subset N$ 是一个坐标邻域；这样便可使用典范坐标来计算。此时
$\theta=\sum p_j\,dq^j$，因而 $\mathcal A_H$ 与经典相空间上的哈密顿作用量泛函（1.3.26）具有完全相同的形式。更形式化地，（1.3.26）写成
\begin{equation}
\begin{split}
 d\mathcal A_H(\gamma)(\xi)
  ={}&\langle\theta(\gamma(1)),\xi(1)\rangle
      -\langle\theta(\gamma(0)),\xi(0)\rangle \\
    &+\int_0^1\omega_0\bigl(\dot\gamma-X_H(\gamma),\xi\bigr)\,dt.
\end{split}
\label{eq:conormal-first-variation}
\end{equation}

\begin{exercise}\label{ex:conormal-first-variation}
证明这个公式。
\end{exercise}

这里回顾一下：在我们的约定中，余切丛上的典范辛形式 $\omega_0$ 由 $\omega_0=-d\theta$ 给出。

在经典相空间中，除了周期边界条件以外，我们还考察过作为自然边界条件的两点边界条件；也就是说，它是使\eqref{eq:conormal-first-variation}中的边界项消失的边界条件。对给定的两个位形点 $x,y\in N$，这相当于把 $\mathcal A_H$ 限制在道路空间
\[
 \mathcal P(T^*N;T_x^*N,T_y^*N)
 =\{\gamma:[0,1]\to T^*N\mid
       \gamma(0)\in T_x^*N,\ \gamma(1)\in T_y^*N\}
\]
上。事实表明，这种自然边界条件容许一个比两点边界条件更一般的类别，我们称之为\emph{余法边界条件}。
\stfhIndex{I083}{yu^fa^bian^jie^tiao^jian}{余法边界条件（conormal boundary condition）}

\begin{definition}[余法丛]\label{def:conormal-bundle}
\stfhIndex{I084}{yu^fa^cong}{余法丛（conormal bundle）}
设 $S\subset N$ 是光滑子流形。$S$ 的\emph{余法丛}定义为
\[
 \nu^*S=\{\alpha\in T_q^*N\mid q\in S,\ \alpha|_{T_qS}=0\}.
\]
\end{definition}

下面的引理给出余法丛的基本性质。它说明余法边界条件确是作用量泛函 $\mathcal A_H$ 的\emph{自然边界条件}，因为这一条件使\eqref{eq:conormal-first-variation}中的边界贡献消失。

\begin{lemma}\label{lem:conormal-tautological-vanishing}
设 $S\subset N$ 是光滑子流形，$i_S:\nu^*S\hookrightarrow T^*N$ 是其余法丛的包含映射。则
$i_S^*\theta=0$。特别地，$i_S^*\omega_0=0$。此外，对任意 $S\subset N$，均有
$\dim\nu^*S=n\;(=\dim N)$。
\end{lemma}

\begin{proof}
设 $\xi\in T_\alpha(\nu^*S)$，则 $d\pi(\xi)\in T_{\pi(\alpha)}S$，所以
\[
 \theta_\alpha(\xi)=\alpha(d\pi(\xi))=0,
\]
因为 $\alpha\in\nu^*_{\pi(\alpha)}S$。
\end{proof}

余法边界条件包含如下两个特例。

\begin{example}\label{ex:conormal-special-cases}
\begin{enumerate}[label=(\arabic*)]
\item 设 $S=\{q\}$ 是 $N$ 中的一点，则 $\nu^*\{q\}=T_q^*N$，即 $T^*N$ 在 $q$ 上的纤维。
\item 当 $S=N$ 时，$\nu^*S=o_N$，即零截面。
\end{enumerate}
\end{example}

现在设 $S_0,S_1\subset N$ 是两个子流形，$\nu^*S_i$（$i=0,1$）是它们的余法丛。考虑 $T^*N$ 上道路组成的子集
\[
 \mathcal L(\nu^*S_0,\nu^*S_1)
 =\{\gamma:[0,1]\to T^*N\mid
       \gamma(0)\in\nu^*S_0,\ \gamma(1)\in\nu^*S_1\}.
\]
用 $d\mathcal A_{(H;S_0,S_1)}$ 表示 $d\mathcal A_H$ 在
$\mathcal L(\nu^*S_0,\nu^*S_1)$ 上的限制，则
\[
 d\mathcal A_{(H;S_0,S_1)}(\gamma)(\xi)
 =\int_0^1\omega\bigl(\dot\gamma-X_H(t,\gamma(t)),\xi(t)\bigr)\,dt.
\]
因此，临界点集 $\operatorname{Crit}(\mathcal A_{(H;S_0,S_1)})$ 与从
$\nu^*S_0$ 到 $\nu^*S_1$ 的哈密顿轨道一一对应；换言之，它与满足
\[
 \dot z=X_H(t,z),\qquad z(0)\in\nu^*S_0,\quad z(1)\in\nu^*S_1
\]
的道路 $z$ 一一对应。

在一般辛流形 $(M,\omega)$ 上，并没有作用量泛函 $\mathcal A_H$ 的直接类似物，因为它的定义本质上使用了 $1$-形式 $\theta$。为了把上述哈密顿方程看成 $M$ 的道路空间上某个泛函的极值道路，我们首先需要发展一般辛流形的大量背景几何。我们从若干基本的辛线性代数开始。

\section{辛线性代数}
\label{sec:symplectic-linear-algebra}

本节紧随温斯坦 (Wn87) 的叙述。

设 $(S,\Omega)$ 是有限维辛向量空间。由
\[
 \widetilde\Omega(v)(w):=\Omega(v,w)
\]
定义的线性映射 $\widetilde\Omega:S\to S^*$ 是斜对称同构。设 $C\subset S$ 是子空间。定义 $C$ 的 $\Omega$-正交补为
\[
 C^\Omega:=\{v\in S\mid \Omega(v,w)=0\quad\forall w\in C\}.
\]

\begin{definition}\label{def:isotropic-coisotropic-lagrangian}
\stfhIndex{I076}{yu^mi^xiang}{余迷向（coisotropic）}
\stfhIndex{I176}{mi^xiang}{迷向（isotropic）}
设 $C\subset S$ 是子空间。
\begin{enumerate}[label=(\arabic*)]
\item 若 $C\subset C^\Omega$，则称 $C$ 为\emph{迷向的}；
\item 若 $C\supset C^\Omega$，则称 $C$ 为\emph{余迷向的}；
\item 若 $C=C^\Omega$，则称 $C$ 为\emph{拉格朗日的}。
\end{enumerate}
\end{definition}

现在说明：当 $C$ 余迷向时，商空间 $C/C^\Omega$ 带有由 $\Omega$ 典范诱导的辛内积。下面这个基本引理是所谓\emph{辛约化}的线性版本。
\stfhIndex{I282}{xin^yue^hua}{辛约化（symplectic reduction）}

\begin{lemma}\label{lem:linear-symplectic-reduction}
设 $(V,\Omega)$ 是辛向量空间，$C\subset V$ 是余迷向子空间，并以 $C^\Omega\subset C$ 表示其零子空间。记 $\iota_C:C\hookrightarrow V$ 与
$\pi_C:C\to C/C^\Omega$ 分别为自然包含与投影。则 $\Omega$ 在商空间 $C/C^\Omega$ 上诱导典范辛双线性形式 $\omega_C$；它由关系
\begin{equation}
 \iota_C^*\Omega=\pi^*\omega_C
 \label{eq:linear-reduction-form}
\end{equation}
所刻画。
\end{lemma}

\begin{proof}
由 $C^\Omega$ 的定义，
\[
 \iota_C^*\Omega|_{C^\Omega}\equiv0,
\]
所以 $\iota_C^*\Omega$ 下降为商空间 $C/C^\Omega$ 上的斜对称 $2$-形式。

还须证明它非退化。设 $v\in C$ 且 $\pi_C(v)\ne0$，即 $v\notin C^\Omega$。我们应找到 $w\in C$ 使 $\Omega(v,w)\ne0$；这由 $C^\Omega$ 的定义立即得到。这就完成了证明。
\end{proof}

\subsection{拉格朗日--格拉斯曼流形}
\label{subsec:lagrangian-grassmannian}
\stfhIndexRangeStart{I182}{la^ge^lang^ri^ge^la^si^man^liu^xing}{拉格朗日--格拉斯曼流形（Lagrangian Grassmannian）}

\begin{definition}\label{def:lagrangian-grassmannian}
以 $\mathcal L(S)$ 表示 $(S,\Omega)$ 的所有拉格朗日子空间组成的集合，以
$\operatorname{Sp}(S,\Omega)\subset GL(S)$ 表示 $(S,\Omega)$ 的自同构所成的子群。
\end{definition}

\begin{proposition}\label{prop:lagrangian-linear-basics}
\begin{enumerate}[label=(\arabic*)]
\item 每个迷向子空间都包含在某个拉格朗日子空间中；
\item $\operatorname{Sp}(S,\Omega)$ 在 $\mathcal L(S)$ 上传递地作用；
\item 每个拉格朗日子空间都有一个拉格朗日补空间。
\end{enumerate}
\end{proposition}

\begin{proof}
把 (1) 与 (2) 留作练习（其证明也可参见 (Ar67)）。这里只证明 (3)。任选 $L$ 的两个补空间 $U_1,U_2$，它们未必是拉格朗日的；亦即
\[
 S=U_1\oplus L=U_2\oplus L.
\]
注意到 $L$ 的补空间集合（记作 $\operatorname{Comp}_L(S)$）是一个以
$\operatorname{Hom}(S/L,L)$ 为模型的仿射空间，我们可以取 $U_1$ 与 $U_2$ 的“差”，记为
\[
 \delta_L(U_1,U_2):S/L\longrightarrow L.
\]
更确切地说，把 $q\in S/L$ 写成 $q=L+v$（其中 $v\in S$）。对 $j=1,2$，$q\cap U_j$ 各有唯一交点 $p_j$。定义
\begin{equation}
 \delta_L(U_1,U_2)(q)=p_2-p_1.
 \label{eq:complement-difference}
\end{equation}
下面把 $\delta_L$ 的若干性质留作练习。
\end{proof}

\begin{exercise}\label{ex:complement-difference-properties}
证明下列各点。
\begin{enumerate}[label=(\arabic*)]
\item $\delta_L(U,U)=0$。
\item $\delta_L$ 满足“余圈条件”
\begin{equation}
 \delta_L(U_1,U_2)+\delta_L(U_2,U_3)=\delta_L(U_1,U_3),
 \label{eq:complement-cocycle}
\end{equation}
特别地，$\delta_L(U_1,U_2)=-\delta_L(U_2,U_1)$。
\item 给定一个补空间 $U_1$ 及映射 $\delta\in\operatorname{Hom}(S/L,L)$，存在唯一补空间 $U_\delta$ 使
$\delta_L(U_1,U_\delta)=\delta$。此时记
\begin{equation}
 U_\delta=U_1+\delta.
 \label{eq:affine-complement-notation}
\end{equation}
\item $U$ 是拉格朗日的，当且仅当 $\delta_L(U,U^\Omega)=0$。
\end{enumerate}
\end{exercise}

下面的一般引理将完成\cref{prop:lagrangian-linear-basics}中陈述 (3) 的证明。

\begin{lemma}\label{lem:lagrangian-complement-average}
设 $L\subset S$ 是 $(S,\Omega)$ 的拉格朗日子空间，$C$ 是任意补空间，未必拉格朗日。则有：
\begin{enumerate}[label=(\arabic*)]
\item $\Omega$-正交补 $C^\Omega$ 也是补空间，并且 $C$ 是拉格朗日的，当且仅当 $C=C^\Omega$。
\item 以 $\iota_L:\operatorname{Comp}_L(S)\to\operatorname{Comp}_L(S)$ 表示由
$\iota_L(C)=C^\Omega$ 定义的对合，则
\[
 \operatorname{Fix}\iota_L=\mathcal L_L(S),
\]
其中 $\mathcal L_L(S)$ 是 $\mathcal L(S)$ 中由 $L$ 的拉格朗日补空间组成的子集。
\item 设 $U_1,U_2\in\operatorname{Comp}_L(S)$，并以
$\delta_L(U_1,U_2)\in\operatorname{Hom}(S/L,L)$ 表示它们的差，则
\begin{equation}
 \delta_L(U_1^\Omega,U_2^\Omega)=\delta_L(U_1,U_2).
 \label{eq:orthogonal-complement-difference}
\end{equation}
\item “平均”$(C+C^\Omega)/2$ 是拉格朗日的。此外，映射
\[
 C\longmapsto\frac{C+C^\Omega}{2};\qquad
 \operatorname{Comp}_L(S)\longrightarrow\mathcal L_L(S)
\]
定义了到拉格朗日补空间的典范投影；而且将其视作
$\operatorname{Comp}_L(S)\to\operatorname{Comp}_L(S)$ 的映射时，这个投影是光滑的。
\end{enumerate}
\end{lemma}

\begin{proof}
因为 $C$ 是 $L$ 的补空间，$C\oplus L=S$，故
\[
 0=(C\oplus L)^\Omega=C^\Omega\cap L^\Omega.
\]
又因 $L$ 是拉格朗日的，$L=L^\Omega$，所以 $C^\Omega\cap L=0$。由
$\dim L=\dim C^\Omega=n$，可知 $C^\Omega\oplus L=S$，这证明了第一项。第二项由此立见。

为证明 (3)，注意：对任意拉格朗日子空间 $L\subset S$ 及每个给定的
$\ell\in L$，内乘 $\ell\mathbin{\lrcorner}\Omega$ 定义了 $S$ 上的线性泛函，其在 $L$ 上的限制为零。因此它定义线性映射
$I_L:L\to(S/L)^*$，并且该映射是同构。这个同构诱导自然等同
\[
 \delta\longmapsto I_L\circ\delta;\qquad
 \operatorname{Hom}(S/L,L)\cong\operatorname{Hom}(S/L,(S/L)^*).
\]
后一个空间恰是 $S/L$ 上二次型组成的空间。

有了这些准备，由 $\Omega$ 的非退化性，陈述 (3) 等价于对所有 $v\in S$ 成立的等式
\begin{equation}
 \Omega(\delta_L(U_1,U_2)v,v)
 =-\Omega(\delta_L(U_1^\Omega,U_2^\Omega)v,v).
 \label{eq:complement-quadratic-sign}
\end{equation}
任取 $v\in S$。对 $i=1,2$，分别以
$p_i\in U_i\cap(L+v)$ 和 $p_i^\Omega\in U_i^\Omega\cap(L+v)$ 表示唯一交点。由定义，
\[
 \delta_L(U_1,U_2)=p_2-p_1,
 \qquad
 \delta_L(U_1^\Omega,U_2^\Omega)=p_2^\Omega-p_1^\Omega.
\]
因此我们需要证明
\begin{equation}
 \Omega(p_2-p_1,v)\,\Omega(p_2^\Omega-p_1^\Omega,v)=0.
 \label{eq:source-complement-product}
\end{equation}
由 $p_i\in U_i$ 和 $p_i^\Omega\in U_i^\Omega$ 的选取，$\Omega(p_i^\Omega,p_i)=0$。又因对所有 $i,j=1,2$ 都有 $p_i,p_j^\Omega\in L+v$，而 $L$ 是拉格朗日的，故
\begin{equation}
 \Omega(p_i^\Omega-v,p_j-v)=0
 \label{eq:lagrangian-affine-coset}
\end{equation}
对所有 $i,j=1,2$ 成立。该式给出
\begin{equation}
 \Omega(p_i^\Omega-p_j,v)=\Omega(p_i^\Omega,p_j)
 \label{eq:orthogonal-intersection-identity}
\end{equation}
对所有 $i,j=1,2$ 成立。现在利用\eqref{eq:orthogonal-intersection-identity}计算
\[
\begin{aligned}
 &\Omega(p_2^\Omega-p_1^\Omega,v)-\Omega(p_2-p_1,v)\\
 &\quad=\Omega(p_2^\Omega-p_1^\Omega-p_2+p_1,v)\\
 &\quad=\Omega(p_2^\Omega-p_2,v)-\Omega(p_1^\Omega-p_1,v)\\
 &\quad=\Omega(p_2^\Omega,p_2)-\Omega(p_1^\Omega,p_1).
\end{aligned}
\]
第三个等号使用了\eqref{eq:orthogonal-intersection-identity}。最后两项消失，因为
$p_i\in U_i$、$p_i^\Omega\in U_i^\Omega$（$i=1,2$）。这证明了\eqref{eq:source-complement-product}，从而证明了陈述 (3)。

最后证明末项。由第一步，$C$ 与 $C^\Omega$ 都在仿射空间
$\operatorname{Comp}_L(S)$ 中，所以平均 $M:=(C+C^\Omega)/2$ 作为
$\operatorname{Comp}_L(S)$ 的元素是有意义的。更确切地说，$C$ 与 $C^\Omega$ 在
$\operatorname{Comp}_L(S)$ 中的平均 $M$ 由方程
\[
 \delta_L(M,C)+\delta_L(M,C^\Omega)=0
\]
唯一确定。现在只须证明 $M$ 是对合
$\iota_L:\operatorname{Comp}_L(S)\to\operatorname{Comp}_L(S)$ 的不动点。另一方面，陈述 (3) 给出
\[
 \delta_L(M^\Omega,C)+\delta_L(M^\Omega,C^\Omega)
 =\delta_L(M,C^\Omega)+\delta_L(M,C)=0.
\]
故 $M^\Omega$ 也是 $C$ 与 $C^\Omega$ 的平均。由平均的唯一性，$M=M^\Omega$；因此由\cref{ex:complement-difference-properties}的陈述 (4)，$M$ 必为拉格朗日的。这就完成了证明。
\end{proof}

这个引理的陈述 (4) 随即证明了\cref{prop:lagrangian-linear-basics}的陈述 (3)。

以上证明立即给出如下推论。

\begin{corollary}\label{cor:symmetric-lagrangian-complement}
设 $L\subset S$ 是拉格朗日子空间，$\delta$ 与 $I_L$ 如上。设 $C\subset S$ 是 $L$ 的一个拉格朗日补空间，$U\in\operatorname{Comp}_L(S)$ 是一个补空间。写
$U=C+\delta$，其中
\[
 \delta\in\operatorname{Hom}(S/L,L)
 \cong\operatorname{Hom}(S/L,(S/L)^*).
\]
则 $U$ 是拉格朗日的，当且仅当 $\delta$ 对称，即 $\delta=\delta^*$。
\end{corollary}

\begin{proof}
在上面的证明中，与 $U$ 相伴的二次型无非是
\[
 \Omega(\delta_L(C,U)v,v),\qquad [v]\in S/L.
\]
式\eqref{eq:complement-quadratic-sign}以及 $C$ 的拉格朗日性蕴含
\[
 \Omega(\delta_L(C,U)v,v)+\Omega(\delta_L(C,U^\Omega)v,v)=0.
\]
利用 $\Omega$ 的斜对称性，可把它改写成
\[
 \Omega(\delta_L(C,U)v,v)=\Omega(v,\delta_L(C,U^\Omega)v).
\]
换言之，作为 $\operatorname{Hom}(S/L,(S/L)^*)$ 中的二次型，
\[
 \delta_L(C,U)^*=\delta_L(C,U^\Omega).
\]
因此，若 $U$ 是拉格朗日的，即 $U=U^\Omega$，则
$\delta_L(C,U)^*=\delta_L(C,U)$，即 $\delta_L(C,U)^*$ 对称。反过来，假设
$\delta_L(C,U)^*$ 对称，则上一等式蕴含
\[
 0=-\delta_L(C,U)+\delta_L(C,U^\Omega)
   =\delta_L(U,C)+\delta_L(C,U^\Omega)
   =\delta_L(U,U^\Omega).
\]
这里后两个等式使用了\cref{ex:complement-difference-properties}中的余圈条件。由该练习的 (3)，方程 $\delta_L(U,U^\Omega)=0$ 蕴含 $U$ 是拉格朗日的。这就完成了证明。
\end{proof}

由此得到下面这个基本定理，它描述了辛向量空间中拉格朗日子空间集合的结构。用
$\mathcal L(S)=\mathcal L(S,\Omega)$ 表示拉格朗日子空间的集合；它通常称为\emph{拉格朗日--格拉斯曼流形}。

\begin{theorem}\label{thm:lagrangian-grassmannian-manifold}
设 $\dim S=2n$。$\mathcal L(S;\Omega)$ 是一个非空的 $n(n+1)/2$ 维流形，以秩为 $n$ 的对称矩阵空间为模型。事实上，对任意给定拉格朗日子空间 $L\subset S$，其拉格朗日补空间的集合典范地成为一个以 $S/L$ 上的对称二次型集合 $S^2(S/L)$ 为模型的仿射空间；特别地，这个集合可缩。
\end{theorem}

\subsection{阿诺尔德分层}
\label{subsec:arnold-stratification}
\stfhSubIndex{I183}{la^ge^lang^ri^ge^la^si^man^liu^xing}{拉格朗日--格拉斯曼流形（Lagrangian Grassmannian）}{a^nuo^er^de^fen^ceng}{阿诺尔德分层（Arnold stratification）}

为此，我们需要对拉格朗日--格拉斯曼流形 $\mathcal L(S,\Omega)$ 的结构作些准备。给定拉格朗日子空间 $L\subset S$，考虑 $\mathcal L(S,\Omega)$ 的分层
\[
 \mathcal L(S,\Omega)\supset\mathcal L_1(S,\Omega;L)\supset\cdots
 \supset\mathcal L_n(S,\Omega;L)=\{L\},
\]
其中，对 $k=1,\ldots,n$，
\[
 \mathcal L_k(S,\Omega;L)
 =\{V\in\mathcal L(S,\Omega)\mid\dim(V\cap L)\ge k\}.
\]
这一分层由阿诺尔德 (Ar67) 引入，是分析拉格朗日--格拉斯曼流形
$\mathcal L(S,\Omega)$ 的基础。我们称它为\emph{关于 $L$ 的阿诺尔德分层}。

由\cref{thm:lagrangian-grassmannian-manifold}，每个
$\mathcal L_k(S,\Omega;L)$ 都是 $\mathcal L(S,\Omega)$ 的闭子集，其开层在
$\mathcal L(S,\Omega)$ 中的余维为 $k(k+1)/2$。特别地，对 $k=1$，阿诺尔德 (Ar67) 证明了关于 $\mathcal L_1(S,\Omega;L)$ 的如下基本事实；我们称之为与 $L\subset S$ 相伴的\emph{马斯洛夫闭链}。
\stfhIndex{I202}{ma^si^luo^fu^bi^lian}{马斯洛夫闭链（Maslov cycle）}

\begin{proposition}\label{prop:maslov-cycle-properties}
设 $L\in\mathcal L(S,\Omega)$，则 $\mathcal L_1(S,\Omega;L)$ 具有如下性质。
\begin{enumerate}[label=(\arabic*)]
\item 它在 $\mathcal L(S,\Omega)$ 中的余维为 $1$。
\item 它具有余定向：存在沿 $\mathcal L_1(S,\Omega;L)$ 的向量场 $\xi$，横截于 $\mathcal L_1(S,\Omega;L)$。
\item 子集 $\mathcal L_1(S,\Omega;L)\setminus\mathcal L_2(S,\Omega;L)$ 是光滑流形，而 $\mathcal L_2(S,\Omega;L)$ 在 $\mathcal L_1(S,\Omega;L)$ 中的余维为 $2$。
\end{enumerate}
\end{proposition}

\begin{proof}
性质 (1) 与 (3) 立即来自维数公式
\[
 \operatorname{codim}\mathcal L_k(S,\Omega;L)=\frac{k(k+1)}2.
\]
为证明 (2)，选取达布基
\[
 \{e_1,\ldots,e_n,f_1,\ldots,f_n\}
\]
使 $\operatorname{span}\{f_1,\ldots,f_n\}=L$，并借辛线性变换
\[
 \Phi:(S,\Omega)\longrightarrow\mathbb C^n;\qquad
 e_j\longmapsto\frac{\partial}{\partial x_j},\quad
 f_k\longmapsto\frac{\partial}{\partial y_k}
\]
把 $L\subset(S,\Omega)$ 等同于 $\sqrt{-1}\mathbb R^n\subset\mathbb C^n$，把
$\operatorname{span}\{e_1,\ldots,e_n\}$ 等同于 $\mathbb R^n$；这里
$z_j=x_j+\sqrt{-1}y_j$ 是 $\mathbb C^n$ 的标准复坐标。

对每个拉格朗日子空间 $\lambda\subset\mathbb C^n$，考虑由
$\theta\mapsto e^{i\theta}\cdot\lambda$ 定义的光滑拉格朗日路径。它在
$\Lambda(n):=\mathcal L(\mathbb R^{2n},\omega_0)$ 上诱导向量场 $\xi$：
\[
 \xi(\lambda)=\left.\frac d{d\theta}\right|_{\theta=0}e^{i\theta}\cdot\lambda.
\]
把 $\xi$ 限制到与 $i\mathbb R^n$ 相伴的马斯洛夫闭链
\[
 \Lambda_1(n)=\mathcal L_1(\mathbb R^{2n},\omega_0;i\mathbb R^n)
 :=\{\lambda\in\Lambda(n)\mid\dim(\lambda\cap i\mathbb R^n)\ge1\}.
\]

设 $\lambda=A\cdot\mathbb R^n\in\Lambda_1(n)\setminus\Lambda_2(n)$，其中 $A$ 是酉矩阵。则 $A\cdot\mathbb R^n\cap i\mathbb R^n$ 的维数为 $1$。取
$0\ne v\in\mathbb R^n$ 使 $Av\in i\mathbb R^n$。将 $A$ 分解为
$A=A_{\mathrm{re}}+iA_{\mathrm{im}}$，关系 $Av\in i\mathbb R^n$ 等价于
\begin{equation}
 A_{\mathrm{re}}v=0,
 \label{eq:real-part-kernel}
\end{equation}
也即 $v\in\ker A_{\mathrm{re}}$。现在考虑
$e^{i\theta}\cdot\lambda=e^{i\theta}A\cdot\mathbb R^n$。

\begin{lemma}\label{lem:rotated-real-part-kernel}
对任意充分小的 $\epsilon>0$，当 $0<|\theta|\le\epsilon$ 时，
\[
 \ker(e^{i\theta}A)_{\mathrm{re}}=\{0\}.
\]
\end{lemma}

\begin{proof}
矩阵的零度是上半连续的，亦即对连续矩阵族而言它不会下降；因此只须证明：对任意满足\eqref{eq:real-part-kernel}的非零向量 $v$，只要 $0<|\theta|$ 充分小，就有
$v\notin\ker(e^{i\theta}A)_{\mathrm{re}}$。注意
\[
 (e^{i\theta}A)_{\mathrm{re}}
 =\cos\theta A_{\mathrm{re}}-\sin\theta A_{\mathrm{im}},\qquad
 (e^{i\theta}A)_{\mathrm{im}}
 =\sin\theta A_{\mathrm{re}}+\cos\theta A_{\mathrm{im}}.
\]
对任意满足\eqref{eq:real-part-kernel}的 $v$，
\[
 (e^{i\theta}A)_{\mathrm{re}}v
 =\cos\theta A_{\mathrm{re}}v-\sin\theta A_{\mathrm{im}}v
 =-\sin\theta A_{\mathrm{im}}v\ne0
\]
对所有 $0<|\theta|<\pi$ 成立。这就完成了证明。
\end{proof}

该引理说明，对所有这样的 $\theta\ne0$，
$e^{i\theta}\cdot\lambda=(e^{i\theta}A)\cdot\mathbb R^n$ 都位于
$\Lambda^0(n):=\Lambda(n)\setminus\Lambda_1(n)$ 中。

特别地，把拉格朗日子空间
$\lambda(\theta):=e^{i\theta}\cdot\lambda=U(\theta)\cdot\mathbb R^n$（其中
$U(\theta)=e^{i\theta}A$）写成
\[
 \lambda(\theta)=\mathbb R^n+iS(\theta)\cdot\mathbb R^n
 =(I+iS(\theta))\cdot\mathbb R^n,
\]
这里 $S(\theta)$ 是某个对称 $n\times n$ 矩阵。此外，
\cref{cor:symmetric-lagrangian-complement}说明，$\lambda'(0)$ 横截于
$\Lambda_1(n)$，当且仅当 $S'(\theta)$ 可逆。

下面计算 $S'(\theta)$。因为
$(U(\theta))_{\mathrm{re}}=(e^{i\theta}A)_{\mathrm{re}}$ 可逆，而且当
$0<|\theta|<\pi$ 时 $U(\theta)_{\mathrm{im}}$ 没有特征值 $-1$，可把 $S(\theta)$ 写成
\[
\begin{aligned}
 S
 &=U_{\mathrm{re}}^{-1}U_{\mathrm{im}}
 =-i(U+\overline U)^{-1}(U-\overline U)\\
 &=-i(I+U^{-1}\overline U)^{-1}(I-U^{-1}\overline U)\\
 &=-i(I+\overline U^{,t}\overline U)^{-1}
          (I-\overline U^{,t}\overline U).
\end{aligned}
\]
令 $D(\theta)=\overline U^{,t}\overline U$，则
\[
 S=-i(I+D)^{-1}(I-D).
\]
注意 $D$ 是只依赖拉格朗日子空间 $\lambda(\theta)$ 的酉对称矩阵。只要
$\theta\ne0$，这个表达式就有定义且光滑依赖于 $\theta$。因为
$D(\theta)=e^{-2i\theta}\overline A^{,t}\overline A$，故
$D'(\theta)=-2iD(\theta)$。由此计算得
\[
 S'(\theta)=-4D(\theta)^2(I+D(\theta))^{-1}
 =-(I+S(\theta)^2).
\]
这是负定矩阵，其特征值均不大于 $-1$。这就证明了
\cref{prop:maslov-cycle-properties}的陈述 (2)。
\end{proof}

特别地，马斯洛夫闭链 $\mathcal L_1(S,\Omega;L)$ 定义了典范上同调类
\[
 \mu=\mu(S,\Omega;L)\in H^1(\mathcal L(S,\Omega),\mathbb Z),
\]
并且它与参考拉格朗日子空间 $L$ 的选取无关。

\begin{exercise}\label{ex:maslov-class-reference-independent}
证明上同调类 $\mu(S,\Omega;L)$ 与 $L$ 的选取无关。
\end{exercise}

上面的证明也说明，切空间 $T_L\mathcal L(S,\Omega)$ 典范同构于
$S/L$ 上对称二次型组成的集合 $S^2(L)$。事实上有如下引理。

\begin{lemma}\label{lem:maslov-tangent-cones}
切锥 $T_L\mathcal L_1(S,\Omega)\subset T_L\mathcal L(S,\Omega)$ 把切空间
$T_L\mathcal L(S,\Omega)\setminus\{0\}$ 分解成 $n+1$ 个互不相交的开锥。
\end{lemma}

\begin{proof}
由\cref{cor:symmetric-lagrangian-complement}，切空间
$T_L\mathcal L(S,\Omega)$ 典范同构于对称二次型集合 $S^2(L)$，而
$T_L\mathcal L(S,\Omega)\setminus T_L\mathcal L_1(S,\Omega)$ 对应于非退化二次型。再者，具有固定符号差的两个非退化二次型可以在相应锥中由连续路径连接。这就完成了证明。
\end{proof}

在同构
$\xi\in T_L\mathcal L(S,\Omega)\longmapsto Q_\xi\in S^2(L)$ 下，有分解
\[
 T_L\mathcal L(S,\Omega)\setminus T_L\mathcal L_1(S,\Omega)
 =\bigsqcup_{k=0}^n C_k(L),
\]
其中 $C_k(L)$ 是开锥
\[
 C_k(L)=\{\xi\in T_L\mathcal L(S,\Omega)\mid Q_\xi\text{ 的指标为 }k\}.
\]
该引理立即给出如下推论。

\begin{corollary}\label{cor:maslov-complement-components}
存在 $L\in\mathcal L(S,\Omega)$ 的邻域 $U$，使得集合
$U\setminus\mathcal L_1(S,\Omega;L)$ 恰有 $n+1$ 个连通分支，并且每个分支的闭包都含 $L$。
\end{corollary}

以 $U_k(L)=U\cap C_k(L)$ 表示与上面的 $C_k(L)$ 相伴的分支。利用这个推论，我们证明如下命题。

\begin{proposition}\label{prop:positive-lagrangian-path-homotopy}
设 $(S,\Omega)$ 是辛向量空间，$L,V_1\subset(S,\Omega)$ 是 $S$ 的两个拉格朗日子空间，且
$L\cap V_1=\{0\}$。考虑满足如下条件的光滑路径
$\alpha:[0,1]\to\mathcal L(S,\Omega)$：
\begin{enumerate}[label=(\arabic*)]
\item $\alpha(0)=L$，$\alpha(1)=V_1$；
\item 对所有 $0<t\le1$，
$\alpha(t)\in\mathcal L_0(S,\Omega;L):=\mathcal L(S,\Omega)\setminus\mathcal L_1(S,\Omega;L)$；
\item $\alpha'(0)\in C_0(L)\subset T_L\mathcal L(S,\Omega)$。
\end{enumerate}
则任意两条这样的路径 $\alpha_1,\alpha_2$ 都相对于端点同伦。
\end{proposition}

\begin{proof}
设 $\alpha_1,\alpha_2:[0,1]\to\mathcal L(S,\Omega)$ 是两条这样的路径。考虑拼接环路
$\alpha:=\alpha_1*\alpha_2^{-1}$：
\[
 \alpha_1*\alpha_2^{-1}=
 \begin{cases}
  \alpha_1(2t),&t\in[0,\tfrac12],\\
  \alpha_2(2(1-t)),&t\in[\tfrac12,1].
 \end{cases}
\]
由条件 (1) 和 (2)，该环路除 $t=0\equiv1$ 外不与马斯洛夫闭链
$\mathcal L_1(S,\Omega;L)$ 相交。另一方面，条件 (3) 蕴含存在某个
$\epsilon_0>0$，使
$\alpha([ -\epsilon_0,0)\cup(0,\epsilon_0])$ 包含于锥形邻域 $U_n(L)$。因此可在
$[-\epsilon,\epsilon]$ 上把 $\alpha$ 从 $L$ 推开并推入
$U_n(L)\subset\mathcal L(S,\Omega)$，同时不改变它在
$S^1\setminus[-2\epsilon,2\epsilon]$ 上的部分。形变后的环路完全不与马斯洛夫闭链
$\mathcal L_1(S,\Omega;L)$ 相交，故可以相对于 $V_1$ 收缩成常值路径 $V_1$。这证明
$\alpha_1$ 与 $\alpha_2$ 相对于端点同伦\footnote{译者注：原书此处误作 $\alpha_0$ 与 $\alpha_1$。}，从而完成证明。
\end{proof}

关于 $\mathcal L(S,\Omega)$ 的拓扑以及 $\mathcal L(S,\Omega)$ 上马斯洛夫指标的描述，可参见 (Ar67)、(GS77) 或 (RS93) 中的进一步优美讨论。后续各章讨论伪全纯圆盘与拉格朗日交弗洛尔理论时，我们还会论及拉格朗日--格拉斯曼流形上的马斯洛夫指标。眼下先回到达布--温斯坦定理的证明。

\section{达布--温斯坦定理}
\label{sec:darboux-weinstein}

考虑一般辛流形上的拉格朗日子流形 $L\subset(M,\omega)$。

\begin{definition}\label{def:lagrangian-immersion-embedding}
\stfhSubIndex{I187}{la^ge^lang^ri^zi^liu^xing}{拉格朗日子流形（Lagrangian submanifold）}{qian^ru}{嵌入（embedding）}
\stfhSubIndex{I188}{la^ge^lang^ri^zi^liu^xing}{拉格朗日子流形（Lagrangian submanifold）}{jin^ru}{浸入（immersion）}
设 $(M,\omega)$ 是 $2n$ 维辛流形，$L$ 是维数为
$n=\frac12\dim M$ 的流形。若浸入（相应地，嵌入）$i:L\to M$ 满足
$i^*\omega=0$，则称它为\emph{拉格朗日浸入}（相应地，\emph{拉格朗日嵌入}）。当 $i$ 是嵌入时，称其像为一个（嵌入的）拉格朗日子流形。
\end{definition}

关于拉格朗日子流形，第一个普遍结果是温斯坦 (Wn73) 的如下邻域定理。

\begin{theorem}[达布--温斯坦定理]\label{thm:darboux-weinstein}
\stfhIndexRangeStart{I093}{da^bu^wen^si^tan^ding^li}{达布--温斯坦定理（Darboux–Weinstein theorem）}
设 $j:L\to(M,\omega)$ 是紧拉格朗日浸入，$Z:L\to T^*L$ 是零截面。则存在从 $T^*L$ 中 $Z$ 的一个邻域出发的浸入
\[
 E:V\longrightarrow M
\]
使得
\begin{enumerate}[label=(\arabic*)]
\item $E^*\omega=\omega_0$；
\item $E\circ Z=j$。
\end{enumerate}
若 $j$ 还是嵌入，则可以把 $E$ 取成嵌入。
\end{theorem}

我们把这个定理约化到温斯坦的形变引理，即引理 2.2.22，以此给出证明。这里需要用到\cref{sec:symplectic-linear-algebra}中发展的基本辛线性代数。以下紧随温斯坦 (Wn87) 的叙述。
\stfhIndexRangeEnd{I182}{la^ge^lang^ri^ge^la^si^man^liu^xing}{拉格朗日--格拉斯曼流形（Lagrangian Grassmannian）}

\subsection{达布--温斯坦定理的证明}
\label{subsec:proof-darboux-weinstein}

先给出证明梗概。
\begin{enumerate}[label=(\arabic*)]
\item \textbf{第一步。}使用上一小节发展的典范辛代数，在一次线性化的层面解决寻找微分同胚的问题。更确切地说，我们要在包含映射 $j:L\to M$ 上找到丛映射
$A:T^*L\to TM$，使丛映射
\[
 dj\oplus A:TL\oplus T^*L\longrightarrow j^*TM
\]
成为 $L$ 上的辛丛映射，亦即在各纤维上为辛映射。再把这个辛映射与典范同构
\[
 T(T^*L)|_{o_L}\cong TL\oplus T^*L
\]
复合；将 $o_L$ 与 $L$ 等同后，它便定义恒等映射之上的辛丛映射
\[
 dj\oplus A:T(T^*L)|_{o_L}\longrightarrow j^*(TM).
\]

\item \textbf{第二步。}选择映射 $E_0:V_0\to M$，其中 $V_0$ 是
$o_L\subset T^*L$ 的一个邻域，并使
\[
 dE_0|_{o_L}=dj\oplus A.
\]
把 $V_0$ 取得充分小，$E_0$ 就成为浸入。

\item \textbf{第三步。}比较 $T^*L$ 上的典范辛形式 $\omega_0$ 与
$V_0\subset T^*L$ 上的 $E_0^*\omega$。由 $E_0$ 的选取，
\begin{equation}
 \omega_0|_{o_L}=E_0^*\omega|_{o_L}.
 \label{eq:darboux-weinstein-forms-agree}
\end{equation}
现在应用形变引理 2.2.22，找到微分同胚 $\psi:V\to V_0$，使
$\psi^*(E_0^*\omega)=\omega_0$ 在 $V$ 上成立，并且
$\psi|_{o_L}=\operatorname{id}_{o_L}$。于是浸入 $E=E_0\circ\psi$ 满足要求。
\end{enumerate}

只须完成第一步。注意 $TL\oplus T^*L$ 与 $j^*TM$ 都是 $L$ 上的辛向量丛。我们要用如下引理。

\begin{lemma}\label{lem:lagrangian-splitting-standard-model}
设 $(S,\Omega)$ 是辛向量空间，并设 $S=L\oplus L'$，其中 $L,L'$ 都是拉格朗日的。由
\[
 B(v)=\widetilde\Omega(v)|_L
\]
定义的映射 $B:L'\to L^*$ 是同构，并且直和映射
\[
 1\oplus B:S=L\oplus L'\longrightarrow L\oplus L^*
\]
是辛同构。
\end{lemma}

\begin{proof}
证明留作练习。
\end{proof}

因为 $j:L\to M$ 是拉格朗日浸入，$TL\subset j^*TM$ 是 $L$ 上的拉格朗日子丛。用 $\Omega$ 表示丛 $j^*TM$ 上逐纤维的辛形式。任选光滑补子丛 $U\subset j^*TM$，亦即满足
$TL\oplus U=j^*TM$ 的 $U$。由\cref{lem:lagrangian-complement-average}，
$U^\Omega$ 是另一个光滑补子丛。再取平均
$C=(U+U^\Omega)/2$，它定义一个拉格朗日子丛。这里 $U^\Omega$ 或 $C$ 的光滑性来自第 2.3.8 节中平均的函子性构造。因此我们得到拉格朗日分裂
$j^*TM=TL\oplus C$。把\cref{lem:lagrangian-splitting-standard-model}用于
$C$ 与 $T^*L$，得到丛同构 $B:C\to T^*L$。现令 $A=B^{-1}$，则
\stfhIndexRangeEnd{I093}{da^bu^wen^si^tan^ding^li}{达布--温斯坦定理（Darboux–Weinstein theorem）}
\[
 1\oplus A:TL\oplus T^*L\longrightarrow TL\oplus C=j^*TM
\]
是辛丛同构。这就完成了第一步，进而完成达布--温斯坦定理的证明。

下面陈述达布--温斯坦定理的带参数版本，并把证明留作练习。

\begin{exercise}[带参数的达布--温斯坦定理]\label{ex:parametrized-darboux-weinstein}
\stfhIndex{I223}{can^shu^hua^da^bu^wen^si^tan^ding^li}{参数化达布--温斯坦定理（parameterized Darboux–Weinstein theorem）}
设 $\psi:[0,1]\times L\to M$ 是拉格朗日同痕，即它是光滑映射，且每个
$\psi_t:L\to M$ 都是拉格朗日嵌入。则存在光滑映射
\[
 \Phi:\bigcup_{t\in[0,1]}V_t\subset[0,1]\times T^*L\longrightarrow M
\]
使得 $\Phi_t:V_t\subset T^*L\to M$ 是辛嵌入，并满足
\begin{enumerate}[label=(\arabic*)]
\item $\Phi_t^*\omega=\omega_0$；
\item 对所有 $t\in[0,1]$，$\Phi_t\circ Z=\psi_t$ 在 $L$ 上成立。
\end{enumerate}
特别地，若 $L$ 紧，则可选择邻域
\[
 \mathcal V\subset\overline{\mathcal V}\subset\bigcap_{t\in[0,1]}V_t,
\]
使得对所有 $t\in[0,1]$，$\Phi_t^*\omega=\omega_0$ 在 $\mathcal V$ 上成立。
\end{exercise}

\section{正合拉格朗日子流形}
\label{sec:exact-lagrangian-submanifolds}

余切丛上的余法丛只是拉格朗日子流形非常特殊的例子；拉格朗日子流形可以定义在一般辛流形上。

在开始一般性研究拉格朗日子流形之前，先看几个例子。

\begin{example}\label{ex:exact-lagrangian-cotangent}
考虑带典范辛形式的余切丛 $M=T^*N$。
\begin{enumerate}[label=(\arabic*)]
\item 对余法丛 $\nu^*S$，有
\begin{equation}
 i^*\theta\equiv0.
 \label{eq:conormal-liouville-zero}
\end{equation}
\item 闭 $1$-形式的图像：设 $\beta$ 是 $N$ 上的闭 $1$-形式，并把 $\beta$ 看作相应的截面映射
\[
 \widetilde\beta:N\longrightarrow T^*N,
\]
其像记作 $\operatorname{Graph}\beta$。在这种情形下，由（2.1.2）可得
\begin{equation}
 d(\widetilde\beta^*\theta)=d\beta=0.
 \label{eq:closed-one-form-graph}
\end{equation}
特别地，光滑函数 $f:N\to\mathbb R$ 的微分 $df$ 的图像是拉格朗日的。
\end{enumerate}
\end{example}

正合 $1$-形式 $df$ 的图像，是所谓正合拉格朗日子流形的特例。

\begin{definition}\label{def:exact-symplectic-manifold}
若辛流形 $(M,\omega)$ 上存在 $1$-形式 $\alpha$ 使 $\omega=d\alpha$，则称它是\emph{正合的}。
\end{definition}

在正合辛流形 $(M,d\alpha)$ 上，$i:L\to M$ 是拉格朗日的，当且仅当
\stfhIndex{I122}{zheng^he^xin^liu^xing}{正合辛流形（exact symplectic manifold）}
$i^*\alpha$ 闭。

\begin{definition}\label{def:exact-lagrangian-submanifold}
正合辛流形 $(M,d\alpha)$ 中的拉格朗日子流形 $i:L\to M$ 称为\emph{正合的}，如果 $i^*\alpha$ 正合。
\stfhIndex{I121}{zheng^he^la^ge^lang^ri^zi^liu^xing}{正合拉格朗日子流形（exact Lagrangian submanifold）}
\end{definition}

正合性定义立即蕴含：存在函数 $f:L\to\mathbb R$ 使 $i^*\alpha=df$。$f$ 的选取在相差加法常数的意义下唯一。稍后会看到，在作用量泛函的变分研究中，把满足 $df=i^*\alpha$ 的配对 $(L,f)$ 看作正合拉格朗日子流形十分重要。

正合拉格朗日子流形的另一个基本性质，是下面这个\emph{周期}消失结论；它是斯托克斯定理的直接推论。

\begin{proposition}\label{prop:exact-lagrangian-period-vanishing}
设 $L\subset(M,\omega)$ 是正合拉格朗日子流形，$\omega=d\alpha$。设
$w:(\Sigma,\partial\Sigma)\to(M,L)$ 是从带边界紧二维曲面 $\Sigma$ 出发的任意光滑映射。则
\[
 \int w^*\omega=0.
\]
\end{proposition}

\begin{proof}
利用 $\omega=d\alpha$ 并应用斯托克斯公式，得到
\[
 \int_\Sigma w^*\omega=\int_\Sigma w^*d\alpha
 =\int_{\partial\Sigma}w^*\alpha.
\]
由 $L$ 的正合性，对 $L$ 上某个函数 $f$ 有 $i_L^*\alpha=df$。因此
$w^*\alpha=d(w^*f)$，故最后一个积分为零。
\end{proof}

\stfhIndex{I078}{wei^xing^kong^jian}{位形空间（configuration space）}
在正合辛流形 $(M,d\alpha)$ 上，可以考虑作用量泛函
\[
 \mathcal A_H(\gamma)=-\int_\gamma\alpha
 -\int_0^1H(t,\gamma(t))\,dt
\]
及其第一变分
\begin{equation}
\begin{split}
 d\mathcal A_H(\gamma)(\xi)
 ={}&\int_0^1\omega\bigl(\xi,\dot\gamma-X_H(t,\gamma)\bigr)\,dt\\
 &+\langle\alpha(\gamma(0)),\xi(\gamma(0))\rangle
  -\langle\alpha(\gamma(1)),\xi(\gamma(1))\rangle.
\end{split}
\label{eq:exact-action-first-variation}
\end{equation}
这无非是把\eqref{eq:conormal-first-variation}中的 $\theta$ 换成 $-\alpha$。周期边界条件仍是使边界贡献消失的自然边界条件。

\begin{question}\label{ques:other-natural-boundary}
这个泛函 $\mathcal A_H$ 是否还有别的自然边界条件？
\end{question}

\begin{remark}\label{rem:d-branes-natural-boundary}
这类问题是物理学家为给具有适当\emph{物理拉格朗日量}的给定系统正确定义\emph{运动方程}而经常提出的问题的原型。在\emph{弦论}中，作为给定系统自然边界条件出现的对象大体对应于\emph{D-膜}的概念。大约自 1996 年物理学家在开弦理论中引入 D-膜概念以来，它一直在当代弦论中起基础作用。
\end{remark}

下面尝试回答上述问题。首先注意，当我们对某些子流形 $L_0,L_1$ 施加边界条件
\[
 \gamma(0)\in L_0,\qquad\gamma(1)\in L_1
\]
时，\eqref{eq:exact-action-first-variation}中的两个边界项分别就是
\[
 i_0^*\alpha(\xi(0)),\qquad i_1^*\alpha(\xi(1)).
\]
现在假设 $L_0,L_1$ 是正合拉格朗日子流形，从而对某些函数
$f_0:L_0\to\mathbb R$ 与 $f_1:L_1\to\mathbb R$ 有
\[
 i_0^*\alpha=df_0,\qquad i_1^*\alpha=df_1.
\]
在此情形下，可以在原作用量泛函 $\mathcal A_H$ 上添加由 $f_0,f_1$ 给出的项，考虑修正后的作用量泛函
$\mathcal A_{H;(f_0,f_1)}:\mathcal P(L_0,L_1)\to\mathbb R$：
\begin{equation}
 \mathcal A_{H;(f_0,f_1)}(\gamma)
 =\mathcal A_H(\gamma)+f_1(\gamma(1))-f_0(\gamma(0)).
 \label{eq:modified-exact-action}
\end{equation}
于是这个修正作用量泛函在 $\mathcal P(L_0,L_1)$ 上的第一变分成为
\[
 d\mathcal A_{H;(f_0,f_1)}(\gamma)(\xi)
 =\int_0^1\omega\bigl(\dot\gamma-X_H(t,\gamma(t))\bigr)\,dt,
\]
从而修正作用量泛函的运动方程为
\[
 \dot z=X_H(z),\qquad z(0)\in L_0,\quad z(1)\in L_1,
\]
这正是我们所要的。由此引出塞德尔 (Se03a) 给出的\emph{正合拉格朗日膜}定义，其中还要给 $L$ 添加一些额外装饰。

\begin{definition}[正合拉格朗日膜]\label{def:exact-lagrangian-brane}
考虑满足 $\omega=d\alpha$ 的 $(M,\omega)$。设
$i:L\hookrightarrow(M,\omega)$ 是正合拉格朗日子流形。若配对 $(L,f)$ 满足
$df=i^*\alpha$，则称它为\emph{正合拉格朗日膜}。
\end{definition}

下面用类似余切丛情形的计算，证明正合拉格朗日膜的一个有趣公式。

\begin{proposition}\label{prop:hamiltonian-flow-exact-brane}
设 $(M,d\alpha)$ 是正合辛流形，$(L,g)$ 是正合拉格朗日膜，包含映射为
$i:L\hookrightarrow M$。哈密顿流 $\phi_H^t$ 诱导一族光滑的正合拉格朗日膜 $(L_t,f_t)$；其中 $L_t=i_t(L)$、$i_t:L\to M$、$f_t:L\to\mathbb R$，并且
\begin{equation}
 i_t=\phi_H^t\circ i,\qquad
 f_t=g+\int_0^t\bigl(H_s+\alpha(X_{H_s})\bigr)\circ i_s\,ds.
 \label{eq:hamiltonian-deformed-brane-phase}
\end{equation}
\end{proposition}

\begin{proof}
由 $f_t$ 的定义，$g=f_0$，所以 $i^*\alpha=df_0$。须证明对所有 $t\in[0,1]$，
\begin{equation}
 i_t^*\alpha=df_t.
 \label{eq:deformed-brane-exactness}
\end{equation}
为此，再次计算
\[
 i_t^*\alpha=(\phi_H^t\circ i)^*\alpha=i^*(\phi_H^t)^*\alpha
\]
的导数。简单计算给出
\begin{equation}
\begin{aligned}
 \frac d{dt}(\phi_H^t)^*\alpha
 &=(\phi_H^t)^*\mathcal L_{X_{H_t}}\alpha\\
 &=(\phi_H^t)^*\bigl(X_{H_t}\mathbin{\lrcorner}d\alpha\bigr)
   +d\bigl((\phi_H^t)^*(X_{H_t}\mathbin{\lrcorner}\alpha)\bigr)\\
 &=d\bigl(H_t\circ\phi_H^t+\alpha(X_{H_t})\circ\phi_H^t\bigr).
\end{aligned}
\label{eq:liouville-form-under-hamiltonian-flow}
\end{equation}
因此，令 $i_0^*\alpha=dg$，便有
\[
\begin{aligned}
 i_t^*\alpha
 &=dg+\int_0^t\frac d{ds}i_s^*\alpha\,ds\\
 &=d\left(g+\int_0^t
   \bigl[H_s\circ\phi_H^s+\alpha(X_{H_s})\circ\phi_H^s\bigr]\circ i\,ds\right)\\
 &=d\left(g+\int_0^t\bigl[H_s+\alpha(X_{H_s})\bigr]\circ i_s\,ds\right).
\end{aligned}
\]
这就由上面对 $f_t$ 的定义证明了\eqref{eq:deformed-brane-exactness}。
\end{proof}

\begin{remark}\label{rem:brane-phase-monodromy}
需要指出，即使 $L_1=L_0$，函数 $g_1$ 也可能不同于 $g_0=g$。正合拉格朗日膜的相位函数所呈现的这种单值性现象，在塞德尔对正合拉格朗日膜的深谷范畴的研究 (Se03a) 中产生了一个有趣结果。关于这一问题的进一步讨论参见 (Se03a)。
\end{remark}

\section{拉格朗日子流形的经典形变}
\label{sec:classical-lagrangian-deformations}
\stfhSubIndexRangeStart{I186}{la^ge^lang^ri^zi^liu^xing}{拉格朗日子流形（Lagrangian submanifold）}{jing^dian^xing^bian}{经典形变（classical deformations）}

达布--温斯坦定理的一个直接推论，完整描述了给定拉格朗日子流形附近的拉格朗日子流形之局部参数化。下面把这一陈述说得更精确。

先引入一个一般定义。

\begin{definition}\label{def:hamiltonian-isotopic-lagrangians}
\stfhSubIndex{I154}{ha^mi^dun}{哈密顿（Hamiltonian）}{ha^mi^dun^tong^hen}{哈密顿同痕（Hamiltonian isotopy）}
若存在含时哈密顿量 $H:[0,1]\times M\to\mathbb R$ 使
\[
 L_1=\phi_H^1(L_0),
\]
则称两个拉格朗日子流形 $L_0,L_1$ \emph{哈密顿同痕}。称同痕
$\{L_t=\phi_H^t(L_0)\}$ 为 $L_0$ 的一个\emph{哈密顿同痕}。
\end{definition}

\subsection{零截面的情形}
\label{subsec:deformations-zero-section}

回顾：一个 $1$-形式（也即余切丛 $T^*M\to M$ 的一个截面）的图像是拉格朗日的，当且仅当这个 $1$-形式闭。此外，任何在 $C^1$ 意义下充分靠近零截面
$o_L\subset T^*L$ 的拉格朗日子流形都是某个 $1$-形式的图像。特别地，如果把 $T^*L$ 的拉格朗日子流形集合
$\operatorname{Lag}(T^*L)$ 看作无限维（弗雷歇）流形，那么零截面的邻域以闭 $1$-形式集合 $Z^1(L)$ 为模型。进一步，任何在 $C^1$ 意义下靠近零截面的\emph{正合}拉格朗日子流形都与零截面哈密顿同痕。这是因为任一这样的拉格朗日子流形都形如
$\operatorname{Graph}(df)$，其中函数 $f$ 在 $C^2$ 意义下靠近常值函数；此时同痕
\[
 t\in[0,1]\longmapsto\operatorname{Graph}(t\,df)
\]
就是所需的哈密顿同痕。

\begin{exercise}\label{ex:zero-section-graph-hamiltonian}
证明上述零截面的同痕由（自治）函数
$H=f\circ\pi:T^*L\to\mathbb R$ 诱导。
\end{exercise}

\begin{exercise}\label{ex:contractible-neighborhood-symp-ham}
证明 $\operatorname{Symp}(M,\omega)$（相应地，$\operatorname{Ham}(M,\omega)$）中恒等映射有一个可缩的 $C^\infty$ 邻域。（提示：考虑辛微分同胚
$\phi:M\to M$ 与其图像
$\operatorname{Graph}\phi\subset M\times M$ 之间的对应；后者关于 $M\times M$ 上的形式
$(-\omega)\oplus\omega$ 是拉格朗日子流形。）
\end{exercise}

更一般地，在余切丛上，拉格朗日子流形模哈密顿同痕有如下局部描述。

\begin{proposition}\label{prop:local-lagrangians-mod-hamiltonian}
在 $C^1$ 意义下靠近零截面的拉格朗日子流形集合，模去 $C^1$-小的哈密顿同痕后，微分同胚于 $H^1(L,\mathbb R)$ 中零点的一个开邻域。
\end{proposition}

\begin{proof}
设 $\mathcal N$ 是 $\operatorname{Lag}(T^*L)$ 中零截面的一个 $C^1$ 邻域。如上所述，对某个 $\delta>0$，可写成
\begin{equation}
 \mathcal N
 =\{\operatorname{Graph}\alpha\mid\alpha\in Z^1(L),\ |\alpha|_{C^1}<\delta\}.
 \label{eq:lagrangian-zero-neighborhood}
\end{equation}
以 $L_\alpha$ 表示与 $\alpha$ 对应的拉格朗日子流形。考虑映射
\[
 L_\alpha\in\mathcal N\longmapsto[\alpha]\in H^1(L,\mathbb R).
\]
显然，该映射满射到 $H^1(L,\mathbb R)$ 中零点的一个开邻域。

假设 $L_\alpha$ 与 $L_{\alpha'}$哈密顿同痕，即对某个函数
$H:[0,1]\times M\to\mathbb R$ 有
$L_{\alpha'}=\phi_H^1(L_\alpha)$，并且对所有 $t\in[0,1]$，
$\phi_H^t(L_\alpha)\subset\mathcal N$。于是可写
\begin{equation}
 \phi_H^t(L_\alpha)=L_{\alpha_t}
 \label{eq:hamiltonian-graph-family}
\end{equation}
其中 $\alpha_t$ 是一族闭 $1$-形式。为计算上同调类 $[\alpha_t]$，对环路
$c:S^1\to L$ 计算 $\int_c\alpha_t$。把 $1$-形式 $\alpha,\alpha_t$ 等同于映射
$\widetilde\alpha,\widetilde\alpha_t:L\to T^*L$，则
\[
 \widetilde\alpha_t=\phi_H^t\circ\widetilde\alpha.
\]
因此，回顾 $\beta=(\widetilde\beta)^*\theta$，有
\[
 \int_c\alpha_t
 =\int_c(\widetilde\alpha_t)^*\theta
 =\int_c(\widetilde\alpha)^*(\phi_H^t)^*\theta
 =\int_{\widetilde\alpha\circ c}(\phi_H^t)^*\theta.
\]
计算
\[
\begin{aligned}
 \int_c\alpha'-\int_c\alpha
 &=\int_0^1\left(\int_c\frac d{dt}\alpha_t\right)dt\\
 &=\int_0^1\left(\int_{\widetilde\alpha\circ c}
     \frac d{dt}(\phi_H^t)^*\theta\right)dt.
\end{aligned}
\]
在上式中把\eqref{eq:liouville-form-under-hamiltonian-flow}用于
$\alpha=-\theta$，得到
\[
 \int_{\widetilde\alpha\circ c}\frac d{dt}(\phi_H^t)^*\theta=0.
\]
故对任意环路 $c$，$\int_c\alpha'-\int_c\alpha=0$。这证明
$[\alpha]=[\alpha']\in H^1(L;\mathbb R)$，因而上述赋值诱导良定义的满射
\[
 \mathcal N/\operatorname{Ham}\longrightarrow H^1(L,\mathbb R),
\]
其像是 $H^1(L,\mathbb R)$ 中零点的一个开邻域。还须证明单射性。设
$[\alpha]=[\alpha']$，亦即 $\alpha'-\alpha=df$。与
\cref{ex:zero-section-graph-hamiltonian}类似，单射性由下面的
\cref{ex:closed-form-graph-hamiltonian}得到。这就完成了证明。
\end{proof}

\begin{exercise}\label{ex:closed-form-graph-hamiltonian}
证明映射
\[
 t\longmapsto\alpha_t:=\alpha+t\,df
\]
以及与它相伴的拉格朗日子流形 $L_{\alpha_t}$，是由函数
$H=f\circ\pi$ 所相伴的哈密顿同痕诱导的。
\end{exercise}

\subsection{一般情形}
\label{subsec:general-lagrangian-deformation}

现在考虑一般的 $(M,\omega)$。设 $L\subset M$ 是拉格朗日嵌入。达布--温斯坦定理给出微分同胚
\stfhIndex{I086}{yu^qie^cong^zuo^biao^ka}{余切丛坐标卡（cotangent bundle chart）}
\[
 \Phi:U\subset M\longrightarrow V\subset T^*L,
\]
它延拓映射 $L\to o_L$，并且 $\Phi^*\omega_0=\omega$。称 $\Phi$ 为 $L$ 的\emph{余切丛坐标卡}。

注意，所有在 $C^1$ 意义下充分靠近给定拉格朗日子流形 $L\subset M$ 的拉格朗日子流形（事实上，这一目的下 $C^0$-接近就够了）都包含在一个给定达布坐标卡中；该坐标卡 $\Phi$ 把 $L$ 映到
$o_L\subset T^*L$。由上面对余切丛的讨论可知，任一这样的、在 $C^1$ 意义下靠近 $L\subset M$ 的拉格朗日子流形，都可由 $L$ 上的闭 $1$-形式参数化；这些形式位于 $Z^1(L)$ 中零点的一个开 $C^1$ 邻域。需要强调，这种参数化依赖余切丛坐标卡，但它的无穷小版本并不依赖。

先研究两个余切丛坐标卡 $\Phi_1,\Phi_2$ 的转移映射；两者均延拓典范等同
$Z:L\to o_L\subset T^*L$。令
\[
 \Phi_{12}=\Phi_2\circ(\Phi_1)^{-1}:V\longrightarrow V'.
\]
它满足 $\Phi_{12}^*\omega_0=\omega_0$ 以及
$\Phi_{12}|_{o_L}=\operatorname{id}|_{o_L}$。回顾：切空间
$T_L\operatorname{Sub}(M)$ 典范同构于 $\Gamma(\nu(L))$，其中
$\nu(L)=TM/TL$ 是 $L$ 的法丛，而 $\operatorname{Sub}(M)$ 是 $M$ 的 $n$ 维子流形集合。

\begin{lemma}\label{lem:normal-cotangent-identification}
有典范同构
\[
 \Gamma(\nu(o_L))\cong\Gamma(T^*L),
\]
它由
$[X]\mapsto(X\mathbin{\lrcorner}\omega_0)|_{To_L}$ 给出，其中 $X$ 是沿
$o_L$ 的向量场。在这个等同下，有自然等同
\begin{equation}
 T_{o_L}\operatorname{Lag}(T^*L)\cong Z^1(L),
 \label{eq:tangent-lag-closed-forms}
\end{equation}
并且 $d\Phi_{12}$ 在 $\Gamma(T^*L)$ 上诱导恒等映射。
\end{lemma}

\begin{proof}
为证明第一项，回顾典范分裂
$T(T^*L)|_{o_L}=TL\oplus T^*L$。先证明上述映射良定义。设 $X,X'$ 是沿
$o_L$ 的两个向量场，且 $[X]=[X']$ 于 $\Gamma(\nu(o_L))$ 中成立。为简便起见，下文只写 $o_L=L$。于是对所有 $x\in L$，
$X(x)-X'(x)\in T_xL$。因为 $L$ 是拉格朗日的，
\[
 (X-X')(x)\mathbin{\lrcorner}\omega_0|_{T_xL}=0,
 \quad\text{即}\quad
 X(x)\mathbin{\lrcorner}\omega_0|_{T_xL}
 =X'(x)\mathbin{\lrcorner}\omega_0|_{T_xL}.
\]
所以 $[X]\mathbin{\lrcorner}\omega_0=[X']\mathbin{\lrcorner}\omega_0$，映射良定义。类似地，它是双射；证明留给读者。

前面已经说明，与 $1$-形式 $\alpha\in\Gamma(T^*L)$ 对应的
$\operatorname{Sub}(T^*L)$ 的切向量，在 $o_L$ 处切于
$\operatorname{Lag}(T^*L)$，当且仅当 $\alpha$ 闭。这就建立了同构\eqref{eq:tangent-lag-closed-forms}。

对 $L$ 上给定的 $1$-形式 $\alpha$，以 $X_\alpha$ 表示沿 $o_L$ 的相伴向量场，它由
$X_\alpha\mathbin{\lrcorner}\omega_0|_{TL}=\alpha$ 定义。则变分
$d\Phi_{12}(X_\alpha)$ 在同构\eqref{eq:tangent-lag-closed-forms}下，对应于
$\Gamma(T^*L)$ 中由
$d\Phi_{12}(X_\alpha)\mathbin{\lrcorner}\omega_0|_{TL}$ 定义的元素。

为求这个截面，对每个给定向量场 $Y\in TL$，直接计算
\[
\begin{aligned}
 d\Phi_{12}(X_\alpha)\mathbin{\lrcorner}\omega_0(Y)
 &=\omega_0(d\Phi_{12}(X_\alpha),Y)\\
 &=\omega_0(X_\alpha,(d\Phi_{12})^{-1}Y)\\
 &=\omega_0(X_\alpha,Y)=\alpha(Y).
\end{aligned}
\]
这就完成证明。这里第二个等式使用 $\Phi_{12}$ 是辛的，第三个等式使用
$\Phi_{12}|_{o_L}=\operatorname{id}|_{o_L}$，最后一个等式使用 $X_\alpha$ 的定义。这就完成了证明。
\end{proof}

用这个引理证明切空间 $T_L\operatorname{Lag}(M,\omega)$ 的如下典范表示。

\begin{proposition}\label{prop:tangent-lag-closed-forms}
设 $\operatorname{Lag}_L(M,\omega)$ 是与 $L$ 具有相同拓扑类型的拉格朗日子流形集合，则有典范同构
\[
 T_L\operatorname{Lag}(M,\omega)\cong Z^1(L).
\]
\end{proposition}

\begin{proof}
因为 $L$ 是拉格朗日的，辛形式 $\omega$ 给出自然同构
\[
 \widetilde\nu(L)\longrightarrow(TL)^*=T^*L;\qquad
 [v]\longmapsto\widetilde\omega(v)|_{TL}.
\]
因此有自然同构
$T_L\operatorname{Sub}(M)\cong\Gamma(T^*L)$。另一方面，给定任意余切丛坐标卡，$L$ 附近的
$\operatorname{Lag}_L(M,\omega)$ 由 $L$ 上的 $C^1$-小闭 $1$-形式集合参数化，故其切空间也由 $Z^1(L)$ 参数化。\Cref{lem:normal-cotangent-identification}说明，切空间的参数化与余切丛坐标卡的选取无关。这就完成了证明。
\end{proof}

切丛
$T\operatorname{Lag}(M,\omega)\to\operatorname{Lag}(M,\omega)$ 含有一个典范分布：在
$L\in\operatorname{Lag}(M,\omega)$ 处，它由正合 $1$-形式集合
$B^1(L)\subset Z^1(L)\cong T_L\operatorname{Lag}(M,\omega)$ 张成。以
$\mathcal E\to\operatorname{Lag}(M,\omega)$ 表示相应子丛。它定义了余秩
$b^1(L)=\dim H^1(L,\mathbb R)$ 的分布。

下面这个在 (Wn90) 中陈述的定理，将在\cref{sec:exact-isotopy-hamiltonian-isotopy}中证明。

\begin{theorem}\label{thm:hamiltonian-leaves}
分布 $\mathcal E\subset T\operatorname{Lag}(M,\omega)$ 是可积子丛；穿过给定拉格朗日子流形
$L\in\operatorname{Lag}(M,\omega)$ 的叶
$\mathcal L_L\subset\operatorname{Lag}(M,\omega)$ 由 $L$ 的哈密顿形变组成。更确切地说，对每个 $L'\in\mathcal L_L$，存在
$\phi\in\operatorname{Ham}(M,\omega)$ 使 $L'=\phi(L)$。事实上，
$\operatorname{Ham}(M,\omega)$ 作用在 $\operatorname{Lag}(M,\omega)$ 上，而
$\mathcal E_L$ 是 $\operatorname{Ham}(M,\omega)$ 在 $L$ 上轨道的切空间。
\end{theorem}

\section{正合拉格朗日同痕 \texorpdfstring{$=$}{=} 哈密顿同痕}
\label{sec:exact-isotopy-hamiltonian-isotopy}

本节依照格罗莫夫 (Gr85)，在任意 $(M,\omega)$ 上引入正合拉格朗日同痕的概念，把它看作 $\operatorname{Lag}(M,\omega)$ 中一种特殊的光滑路径；并证明，这个概念等价于由周围流形 $(M,\omega)$ 上的整体哈密顿流诱导的（嵌入）拉格朗日子流形同痕。我们把后一种同痕称为\emph{整体哈密顿同痕}。对一般拉格朗日同痕而言，这种延拓性质不成立；这是两类同痕之间的一个本质区别。作为推论，本节末尾将证明\cref{thm:hamiltonian-leaves}。

设 $\psi:[0,1]\times L\to M$ 是一族嵌入 $\psi_t:L\to M$ 的同痕。拉回 $\psi^*\omega$ 分解成
\begin{equation}
 \psi^*\omega=dt\wedge\alpha+\beta,
 \label{eq:isotopy-form-decomposition}
\end{equation}
其中 $\alpha,\beta$ 是 $L$ 上的形式，满足
\[
 \frac\partial{\partial t}\mathbin{\lrcorner}\alpha=0
 =\frac\partial{\partial t}\mathbin{\lrcorner}\beta.
\]
事实上，完全用 $\psi^*\omega$ 可写成
\[
 \alpha=\frac\partial{\partial t}\mathbin{\lrcorner}\psi^*\omega,
 \qquad
 \beta=\psi^*\omega-dt\wedge\alpha.
\]
可以把 $\alpha,\beta$ 看作 $L$ 上随 $t$ 变化的形式族。

\begin{definition}\label{def:lagrangian-isotopy}
若在上述分解中 $d\beta=0$，则称同痕
$\psi:[0,1]\times L\to M$ 为\emph{拉格朗日同痕}。
\end{definition}

\begin{proposition}\label{prop:lagrangian-isotopy-closed-form}
以 $i_t:L\to[0,1]\times L$ 表示由 $i_t(x)=(t,x)$ 定义的显然包含。则同痕 $\psi$ 是拉格朗日的，当且仅当对所有 $t\in[0,1]$，$i_t^*\alpha$ 是 $L$ 上的闭 $1$-形式。
\end{proposition}

\begin{proof}
假设 $d\beta=0$。因 $\omega$ 闭，对\eqref{eq:isotopy-form-decomposition}取外微分得到 $0=dt\wedge d\alpha$。再与 $\partial/\partial t$ 作内乘，得到
\[
 d\alpha-dt\wedge\left(\frac\partial{\partial t}
     \mathbin{\lrcorner}d\alpha\right)=0.
\]
用 $i_t$ 拉回该式，得到 $d(i_t^*\alpha)=0$。

反过来，假设 $d(i_t^*\alpha)=0$，则 $i_t^*(d\alpha)=0$。换言之，只要
$X,Y$ 切于纤维 $[0,1]\times L\to[0,1]$，就有 $d\alpha(X,Y)=0$。因此 $3$-形式 $dt\wedge d\alpha$ 在 $[0,1]\times L$ 上必为零。再次对\eqref{eq:isotopy-form-decomposition}取外微分，并使用 $\psi^*\omega$ 的闭性，便得到 $d\beta=0$。这就完成了证明。
\end{proof}

这引出如下定义。

\begin{definition}\label{def:exact-lagrangian-isotopy}
若拉格朗日同痕 $\psi$ 的 $i_t^*\alpha$ 对所有 $t\in[0,1]$ 都正合，则称 $\psi$ \emph{正合}。称这样的同痕为\emph{正合拉格朗日同痕}。
\end{definition}

\begin{lemma}\label{lem:exact-isotopy-primitive}
若 $\psi$ 是正合拉格朗日同痕，则存在函数
$h:[0,1]\times L\to\mathbb R$ 使
\[
 \psi^*\omega=dh\wedge dt.
\]
此外，若 $h'$ 是另一个这样的函数，则 $h'-h$ 只依赖于 $t$，而不依赖于 $L$。
\end{lemma}

\begin{proof}
由定义，对每个 $t\in[0,1]$ 都存在 $h_t:L\to\mathbb R$ 使
$i_t^*\alpha=dh_t$。这个选取在相差常数 $c(t)$ 的意义下唯一。回顾
\cref{def:lagrangian-isotopy}之前的说明，形式 $\alpha,\beta$ 由 $\psi$ 唯一确定。这证明了引理的最后一项。

唯一还须保证的是 $h$ 光滑依赖于 $t$。在 $L$ 的每个连通分支中固定一点
$x_0$，要求对所有 $t\in[0,1]$ 都有 $h_t(x_0)=0$，即可做到这一点。然后定义
$h:[0,1]\times L\to\mathbb R$ 为 $h(t,x)=h_t(x)$。
\end{proof}

\begin{exercise}\label{ex:smooth-primitive-exact-isotopy}
补全上面的证明。更确切地说，证明上述函数 $h$ 确实光滑，并且满足
$\psi^*\omega=dh\wedge dt$。
\end{exercise}

\begin{remark}\label{rem:exact-isotopy-immersions}
需要指出，上述拉格朗日同痕和正合拉格朗日同痕的定义也可用于浸入，而不只用于嵌入。
\end{remark}

下面的延拓定理把\emph{嵌入的正合拉格朗日同痕}同一般拉格朗日同痕以及浸入的正合拉格朗日同痕区别开来。由于这个定理，对嵌入而言，我们经常不区分“哈密顿同痕”和“正合拉格朗日同痕”这两个术语。

\begin{theorem}\label{thm:exact-lagrangian-isotopy-extension}
对任意由拉格朗日嵌入组成的正合拉格朗日同痕
$\psi:[0,1]\times L\to M$，存在 $H:[0,1]\times M\to\mathbb R$ 使
\begin{equation}
 \psi(t,x)=\phi_H^t(\psi(0,x))=\phi_H^t(\psi_0(x)).
 \label{eq:exact-isotopy-ambient-hamiltonian}
\end{equation}
\end{theorem}

分两步证明。第一步，找函数 $H':[0,1]\times M\to\mathbb R$ 使
\begin{equation}
 \frac{\partial\psi}{\partial t}
 \equiv X_{H'}(\psi(t,x))|_{L_t}\pmod{TL_t},
 \label{eq:exact-isotopy-normal-extension}
\end{equation}
其中 $L_t=\psi_t(L)$ 是 $\psi_t$ 的像。第二步，修改 $H'$ 得到 $H$，使
\begin{equation}
 \frac{\partial\psi}{\partial t}=X_H(\psi(t,x))
 \label{eq:exact-isotopy-full-extension}
\end{equation}
对 $(t,x)\in[0,1]\times L$ 成立。先从一个一般引理开始。

\begin{lemma}\label{lem:tangent-hamiltonian-extension}
设 $L$ 是 $n$ 维流形。
\begin{enumerate}[label=(\arabic*)]
\item 若 $g:L\to(M,\omega)$ 是拉格朗日嵌入，$\xi\in\Gamma(TL)$ 是 $L$ 上的向量场，则存在函数 $f:M\to\mathbb R$ 使
\begin{equation}
 X_f(g(x))=dg(\xi(x)),\qquad x\in L.
 \label{eq:tangent-vector-field-hamiltonian-extension}
\end{equation}
\item 对给定的、由嵌入组成的拉格朗日同痕
$\psi:[0,1]\times L\to M$，以及 $L$ 上一族光滑向量场
$\xi=\{\xi_t\}_{0\le t\le1}$，存在光滑函数
$F:[0,1]\times M\to\mathbb R$，使对所有 $t\in[0,1]$，
\begin{equation}
 X_{F_t}(\psi_t(x))=d\psi_t(\xi_t(x)).
 \label{eq:parametrized-tangent-hamiltonian-extension}
\end{equation}
\end{enumerate}
\end{lemma}

\begin{proof}
乘上一个截断函数后，达布--温斯坦定理把 (1) 的证明约化到
$L=o_L$、$M=T^*L$ 带典范辛形式 $\omega_0$ 的情形。但若 $\xi$ 是 $o_L$ 上的向量场，可以显式选取函数 $f_\xi:T^*L\to\mathbb R$：
\begin{equation}
 f_\xi(\alpha)=\langle\alpha,\xi(\pi(\alpha))\rangle.
 \label{eq:fiber-linear-hamiltonian}
\end{equation}
下面的练习将完成这一断言的证明。

对 (2)，使用带参数的达布--温斯坦定理
（\cref{ex:parametrized-darboux-weinstein}），并定义
$F_\xi:[0,1]\times T^*L\to\mathbb R$ 为
\[
 F(t,\alpha)=\langle\alpha,\xi_t(\pi(\alpha))\rangle.
\]
这就完成证明。
\end{proof}

\begin{exercise}\label{ex:fiber-linear-hamiltonian}
证明在 $L$ 上 $X_{f_\xi}\circ g=\xi$。
\end{exercise}

另一个重要引理是下面的一般延拓引理。

\begin{lemma}[延拓引理]\label{lem:function-extension}
设 $\pi:E\to N$ 是向量丛，$h:N\to\mathbb R$ 是 $C^1$ 函数。则存在
$C^1$ 函数 $H':E\to\mathbb R$，使
\[
 H'|_N=h,\qquad dH'|_{TN}=d(h\circ\pi)|_{TN},
\]
并且 $H'$ 的支集可以放进 $N$ 的一个邻域 $U$ 中，而该邻域可以取得任意小。
\end{lemma}

\begin{proof}
设 $r_t:E\to E$ 是 $E$ 到零截面 $o_E\cong N$ 的收缩：对 $e\in E$，
\[
 r_t(e)=
 \begin{cases}
  te,&0\le t\le1,\\
  \pi(e),&t=0.
 \end{cases}
\]
固定映射 $\chi:[0,1]\to[0,1]$，使得 $\chi(t)$ 在 $t=0$ 附近恒为 $0$、
$\chi(1)=1$ 且 $\chi'(t)\ge0$。

若定义含时向量场
\[
 X_t=
 \begin{cases}
  \dfrac{dr_{\chi(t)}}{dt}\circ r_{\chi(t)}^{-1},&0<t\le1,\\[4pt]
  0,&t=0,
 \end{cases}
\]
则对 $0\le t\le1$，
\[
 \frac d{dt}r_{\chi(t)}^*\alpha
 =r_{\chi(t)}^*\bigl(d(X_t\mathbin{\lrcorner}\alpha)
       +X_t\mathbin{\lrcorner}d\alpha\bigr).
\]
当形式 $\alpha$ 闭时，有
\[
 \frac d{dt}r_{\chi(t)}^*\alpha
 =r_{\chi(t)}^*\bigl(d(X_t\mathbin{\lrcorner}\alpha)\bigr).
\]
于是
\[
\begin{aligned}
 \alpha-r_0^*\alpha
 &=\int_0^1r_{\chi(t)}^*\bigl(d(X_t\mathbin{\lrcorner}\alpha)\bigr)\,dt\\
 &=d\int_0^1r_{\chi(t)}^*(X_t\mathbin{\lrcorner}\alpha)\,dt,
\end{aligned}
\]
从而对闭 $1$-形式 $\alpha$，
\[
 \alpha=\pi^*\alpha+d\int_0^1
 r_{\chi(t)}^*(X_t\mathbin{\lrcorner}\alpha)\,dt.
\]
把它用于形式 $\alpha=d(\rho h\circ\pi)$，其中 $h\circ\pi:E\to\mathbb R$，而
$\rho:E\to\mathbb R$ 满足
\[
 \rho=
 \begin{cases}
  1,&\text{在 $N$ 的一个邻域 $U$ 上},\\
  0,&\text{远离 $N$ 时}.
 \end{cases}
\]
现在考虑函数 $H':E\to\mathbb R$：
\[
 H'=\rho\left(h\circ\pi+
 \int_0^1r_{\chi(t)}^*
   \bigl(X_t\mathbin{\lrcorner}d(\rho\circ\pi)\bigr)\,dt\right).
\]
因为对所有 $t\in[0,1]$ 都有
$r_{\chi(t)}|_N=\operatorname{id}|_N$ 且 $X_t\equiv0$ 于 $N$ 上，所以
$H'|_N=h$，并且 $dH'|_{TN}=d(h\circ\pi)|_{TN}$。这就完成了证明。
\end{proof}

\begin{exercise}\label{ex:parametrized-extension-lemma}
陈述并证明延拓引理的带参数版本。
\end{exercise}

利用这些引理，证明\cref{thm:exact-lagrangian-isotopy-extension}。

\begin{proof}[\cref{thm:exact-lagrangian-isotopy-extension}的证明]
注意 $\partial\psi/\partial t$ 是沿 $\psi$ 的向量场，即
$(\partial\psi/\partial t)(x)\in T_{\psi(x)}M$。由
\cref{lem:normal-cotangent-identification}，\eqref{eq:exact-isotopy-normal-extension}等价于
\begin{equation}
 \frac{\partial\psi}{\partial t}\mathbin{\lrcorner}\omega
 =dH'(\psi(t,x))|_{L_t}\in\Gamma(T^*L_t)
 \label{eq:exact-isotopy-one-form-extension}
\end{equation}
对 $t\in[0,1]$ 成立。另一方面，因为 $\psi$ 是正合拉格朗日同痕，对某个
$h':[0,1]\times L\to\mathbb R$ 有
\[
 \psi^*\omega=dh'\wedge dt,
\]
亦即
\[
 dh_t'=\frac\partial{\partial t}\mathbin{\lrcorner}\psi^*\omega
 =\omega\left(d\psi_t\!\left(\frac\partial{\partial t}\right),
 d\psi_t(\,\cdot\,)\right).
\]
因此\eqref{eq:exact-isotopy-one-form-extension}化为
\[
 d(h_t'\circ\psi_t^{-1})|_{TL_t}
 =\frac{\partial\psi}{\partial t}\mathbin{\lrcorner}\omega.
\]
（\emph{这里本质地使用了映射 $\psi_t$ 的嵌入性。}）对 $y\in\psi_t(L)$ 令
$h(t,y)=h_t'\circ\psi_t^{-1}(y)$，则\eqref{eq:exact-isotopy-one-form-extension}化为上面练习中的带参数延拓引理。这完成了证明的第一步：找到了 $M$ 上满足\eqref{eq:exact-isotopy-normal-extension}的 $H'$。

于是向量场
\[
 \frac{\partial\psi}{\partial t}(x)-X_{H'}(t,\psi(t,x))
\]
切于 $\psi_t(L)$。因此由\cref{lem:tangent-hamiltonian-extension}，可找到
$F:[0,1]\times M\to\mathbb R$，使
\[
 \frac{\partial\psi}{\partial t}(x)-X_{H'}(t,\psi(t,x))
 =X_F(t,\psi(t,x))
\]
对 $(t,x)\in[0,1]\times L$ 成立。取 $H=H'+F$，定理得证。
\end{proof}

现在可以证明温斯坦的\cref{thm:hamiltonian-leaves}。由
\cref{thm:exact-lagrangian-isotopy-extension}，只须证明对任意
$L\in\operatorname{Lag}(M,\omega)$，
\[
 B^1(L)=\operatorname{span}\{X_f\mathbin{\lrcorner}\omega|_{TL}
       \mid f:M\to\mathbb R\}
 =T_L(\operatorname{Ham}(M,\omega)\cdot L).
\]
但这正是\cref{thm:exact-lagrangian-isotopy-extension}的直接推论，故得证。
\stfhSubIndexRangeEnd{I186}{la^ge^lang^ri^zi^liu^xing}{拉格朗日子流形（Lagrangian submanifold）}{jing^dian^xing^bian}{经典形变（classical deformations）}

\section{拉格朗日子流形的构造}
\label{sec:construction-lagrangian-submanifolds}

\subsection{拉格朗日手术}
\label{subsec:lagrangian-surgery}
\stfhIndex{I189}{la^ge^lang^ri^shou^shu}{拉格朗日手术（Lagrangian surgery）}

设 $L_1,L_2$ 是 $(M,\omega)$ 中一对已定向拉格朗日子流形，在 $p_{12}$ 处横截相交。固定该配对的次序为 $(L_1,L_2)$。总可以在 $p_{12}$ 的邻域 $U$ 中选取达布坐标卡
$I:U\to V\subset\mathbb C^n$，使 $I(p_{12})=0$，并且
\[
 I(L_1\cap U)=\mathbb R^n\cap V,\qquad
 I(L_2\cap U)=\sqrt{-1}\mathbb R^n\cap V.
\]
下面会用超出实际所需的细节解释这一构造；希望这种具体描述对于更几何地研究拉格朗日手术有所帮助，尤其是研究手术下全纯多边形模空间的蜕变时。这样的应用参见 (FOOO07)。

证明来自达布定理的一个版本，但强烈依赖下面这个著名的辛线性代数事实。

\begin{lemma}\label{lem:sp-transitive-transverse-lagrangian-pairs}
线性辛群 $\operatorname{Sp}(2n)$ 在拉格朗日子空间横截配对的集合上传递地作用。
\end{lemma}

\begin{exercise}\label{ex:sp-transitive-transverse-lagrangian-pairs}
证明这个引理。
\end{exercise}

需要指出，$U(n)\subset\operatorname{Sp}(2n)$ 在这类配对的集合上并不传递地作用。例如，酉群 $U(n)$ 的作用保持 $\mathbb C^n$ 的埃尔米特结构，因而保持两个拉格朗日子空间之间的“角度”。

设 $\epsilon$ 是充分接近 $0$ 的实数。选取函数
$f_\epsilon:\mathbb R^n\setminus\{0\}\to\mathbb R$：
\begin{equation}
 f_\epsilon(x)=\epsilon\log|x|,
 \label{eq:surgery-log-function}
\end{equation}
并把 $\operatorname{Graph}df_\epsilon$ 记作
$H_\epsilon\subset\mathbb C^n$。它是
$T^*(\mathbb R^n\setminus\{0\})\cong
 \mathbb C^n\setminus\sqrt{-1}\mathbb R^n$ 中的拉格朗日子流形；当
$|x|\to0$ 时渐近于 $\sqrt{-1}\mathbb R^n$，当 $|x|\to\infty$ 时渐近于
$\mathbb R^n$。注意
\[
 df_\epsilon(x)=\epsilon\frac{x\cdot dx}{|x|^2}
 =\frac\epsilon{|x|^2}\sum_{j=1}^n x_j\,dx_j,
\]
所以在坐标中可写
\begin{equation}
 H_\epsilon=
 \left\{(z_1,\ldots,z_n)\ \middle|\
 y_j=\frac{\epsilon x_j}{|x|^2},\ j=1,\ldots,n\right\}.
 \label{eq:surgery-handle-model}
\end{equation}
这里以 $z_j=x_j+\sqrt{-1}y_j$（$j=1,\ldots,n$）表示 $\mathbb C^n$ 的复坐标。

令 $\tau:\mathbb C^n\to\mathbb C^n$ 为关于对角线
\[
 \Delta=\{(z_1,z_2,\ldots,z_n)\mid x_i=y_i\}
\]
的反射，即映射
\[
 (x_1+\sqrt{-1}y_1,\ldots,x_n+\sqrt{-1}y_n)
 \longmapsto
 (y_1+\sqrt{-1}x_1,\ldots,y_n+\sqrt{-1}x_n).
\]
式\eqref{eq:surgery-handle-model}蕴含 $|x|^2|y|^2=|\epsilon|^2$，所以
\[
 H_\epsilon=
 \left\{(z_1,\ldots,z_n)\ \middle|\
 x_j=\frac{\epsilon y_j}{|y|^2},\ j=1,\ldots,n\right\}.
\]
换言之，$\tau(H_\epsilon)=H_\epsilon$。容易验证
\[
 \inf\{|\vec z|\mid\vec z\in H_\epsilon\}=\sqrt{2|\epsilon|},
\]
并且最小值在 $(n-1)$-球面
\[
 \{(x,y)\mid |x|=|y|=\sqrt\epsilon\}\subset H_\epsilon
\]
上的点处取得。

接着考虑函数 $\rho:\mathbb R_+\to\mathbb R$，使
\[
 \rho=
 \begin{cases}
  \log r-|\epsilon|,&r\le\sqrt{|\epsilon|}S_0,\\
  \log\sqrt{|\epsilon|}S_0,&r\ge2\sqrt{|\epsilon|}S_0,
 \end{cases}
 \qquad
 \rho'(r)\ge0,\quad\rho''(r)\le0,
\]
其中 $S_0$ 是充分大的数，而 $\epsilon$ 取得使
$\sqrt{|\epsilon|}S_0$ 充分小。定义函数
$\widetilde f_\epsilon:\mathbb R^n\setminus\{0\}\to\mathbb R$：
\begin{equation}
 \widetilde f_\epsilon(x)=\epsilon\rho(|x|).
 \label{eq:cutoff-surgery-function}
\end{equation}

考虑图像 $\operatorname{Graph}d\widetilde f_\epsilon$，并定义拉格朗日子流形 $H_\epsilon'$ 使其满足：
\begin{enumerate}[label=(\arabic*)]
\item $\tau(H_\epsilon')=H_\epsilon'$；
\item
\[
\begin{aligned}
 &\{(x_1+\sqrt{-1}y_1,\ldots,x_n+\sqrt{-1}y_n)
       \mid \forall i\;x_i\ge y_i\}\cap H_\epsilon'\\
 &\quad=
 \{(x_1+\sqrt{-1}y_1,\ldots,x_n+\sqrt{-1}y_n)
       \mid \forall i\;x_i\ge y_i\}
       \cap\operatorname{Graph}d\widetilde f_\epsilon.
\end{aligned}
\]
\end{enumerate}
由构造，$H_\epsilon'$ 在 $\tau$ 下不变，并且在
$I(U)\subset\mathbb R^{2n}$ 中以 $0$ 为中心的球
$B^{2n}(2\sqrt{|\epsilon|}S_0)$ 外，
$H_\epsilon'=\mathbb R^n\cup\sqrt{-1}\mathbb R^n$。因此，对给定有序配对
$(L_1,L_2)$，可以构造拉格朗日子流形 $L_\epsilon\subset M$，使
\[
 L_\epsilon-U=L_1\cup L_2-U,\qquad
 I(L_\epsilon\cap U)=H_\epsilon'\cap V.
\]

\begin{definition}\label{def:lagrangian-connected-sum}
对给定有序配对 $(L_1,L_2)$，称 $L_\epsilon$ 为 $L_1$ 与 $L_2$ 在
$p_{12}\in L_1\cap L_2$ 处的\emph{拉格朗日连通和}，并写作
\stfhIndex{I181}{la^ge^lang^ri^lian^tong^he}{拉格朗日连通和（Lagrangian connected sum）}
$L_1\#_\epsilon L_2=L_\epsilon$。
\end{definition}

注意，如果在 $p\in L_1\cap L_2$ 处交换配对 $(L_1,L_2)$ 的次序，并同时改变
$\epsilon$ 的符号，则所得拉格朗日子流形同构。事实上，
\[
 \mathbb R^n\#_\epsilon\sqrt{-1}\mathbb R^n
 =\sqrt{-1}\mathbb R^n\#_{-\epsilon}\mathbb R^n.
\]

称原像
\[
 (L_1\#_\epsilon L_2)\cap U=I^{-1}(H_\epsilon'\cap V)
\]
为\emph{拉格朗日手柄}，称其子午球面 $S^{n-1}$ 为拉格朗日手术
\stfhIndex{I184}{la^ge^lang^ri^shou^bing}{拉格朗日手柄（Lagrangian handle）}
$L_\epsilon$ 的\emph{消失闭链}。
\stfhIndex{I299}{xiao^shi^bi^lian}{消失闭链（vanishing cycle）}

如果 $\epsilon>0$，并且作为已定向向量空间有
\[
 T_{p_{12}}L_1\oplus T_{p_{12}}L_2
 =(-1)^{n(n-1)/2+1}T_{p_{12}}M,
\]
则称配对 $L_1,L_2$ 或其相伴拉格朗日手术 $L_1\#_\epsilon L_2$ 在
\stfhSubIndex{I190}{la^ge^lang^ri^shou^shu}{拉格朗日手术（Lagrangian surgery）}{zheng^de}{正的（positive）}
$p_{12}$ 处为\emph{正的}；否则称为\emph{负的}。（这里给 $T_{p_{12}}M$ 配备辛定向。）例如，取
$L_1=\mathbb R$、$L_2=\sqrt{-1}\mathbb R\subset\mathbb C$，并在
$\mathbb R,\sqrt{-1}\mathbb R$ 上取标准定向，则当 $\epsilon>0$ 时，拉格朗日手术
$L_1\#_\epsilon L_2$ 是负的。（这一例子可直接推广到
$L_1=\mathbb R^n$、$L_2=\sqrt{-1}\mathbb R^n\subset\mathbb C^n$。）容易验证，只有正手术才容许把 $L_1,L_2$ 上的定向粘合起来，使手术
$L_1\#_\epsilon L_2$ 带有相容定向。（当
$L_1=\mathbb R$、$L_2=\sqrt{-1}\mathbb R\subset\mathbb C$ 时，显然不可能给
$L_1\#_\epsilon L_2$ 一个同时与两个 $L_i$ 的标准定向相容的定向；对
$L_1=\mathbb R^n$、$L_2=\sqrt{-1}\mathbb R^n\subset\mathbb C^n$ 也有类似结论。）

一般而言，即使 $L_1\#_\epsilon L_2$ 与 $L_2\#_\epsilon L_1$ 都可定向，前者甚至也不同痕于后者（或 $L_1\#_{-\epsilon}L_2$）。另一方面，容易验证：对同号的不同 $\epsilon$，各
$L_1\#_\epsilon L_2$ 彼此拉格朗日同痕。不过一般而言，在固定边界的条件下，它们彼此并不哈密顿同痕。

因为上述构造局限在交点的一个邻域内，也可以在拉格朗日浸入的\emph{自交}点处实施手术。设 $L\subset M$ 是浸入拉格朗日子流形，在
$p\in L\subset M$ 处有横截自交。取 $M$ 中 $p$ 的小邻域 $U$，考察
$U\cap L$；它在 $p$ 处有两个分支。为实施拉格朗日手术，需要给这两个分支排序。沿用 (Po91b) 的术语，引入如下定义。

\begin{definition}\label{def:surgery-equipment}
\stfhIndex{I118}{pei^bei}{配备（equipment）}
设 $L,p$ 如上。$p$ 处的一个\emph{配备}，是 $L$ 在 $p$ 处两个分支的一个次序；等价地，它是相应两个切空间在 $p$ 处的一个次序。
\end{definition}

选定 $L$ 的一个自交点处的配备，就可以实施拉格朗日手术并减少交点数。这样，可以从一个拉格朗日浸入得到拉格朗日嵌入。由格罗莫夫的
$h$-原理，拉格朗日浸入存在性的障碍纯属拓扑；相比之下，拉格朗日嵌入的存在性有更精细的硬障碍（此类定理参见 (Gr85)）。遗憾的是，由拉格朗日手术得到的嵌入往往具有复杂的拓扑。

\begin{example}\label{ex:whitney-sphere-surgery}
\stfhIndex{I305}{hui^te^ni^qiu^mian}{惠特尼球面（Whitney sphere）}
格罗莫夫 (Gr85) 证明，$\mathbb C^n$ 中任意紧嵌入拉格朗日子流形 $L$ 必须满足
$H_1(L;\mathbb R)\ne0$。特别地，$\mathbb C^n$ 中不存在嵌入拉格朗日球面
$S^n$。但 $\mathbb C^n$ 中有一个著名的浸入球面，称为惠特尼球面，并且可以使它成为拉格朗日的。事实上，浸入
$\psi:S^n\to\mathbb C^n$ 有如下显式公式：
\begin{equation}
 \psi(x_1,\ldots,x_{n+1})
 =\frac1{1+x_{n+1}^2}
 \bigl(x_1(1+ix_{n+1}),x_2(1+ix_{n+1}),\ldots,x_n(1+ix_{n+1})\bigr).
 \label{eq:whitney-sphere-immersion}
\end{equation}
它在 $\mathbb C^n$ 的原点 $0$ 处有唯一一个横截自交点，其原像是
$S^n\subset\mathbb R^{n+1}$ 的南北两极。也容易验证 $\psi$ 是拉格朗日的。在此交点对它实施拉格朗日手术，就得到
$S^1\times S^{n-1}$ 到 $\mathbb C^n$ 中的拉格朗日嵌入。

$n=2$ 的情形尤其有趣。以 $(z_1,z_2)$ 表示 $\mathbb C^2$ 的复坐标。这时公式化为
\begin{equation}
 \psi(x,y,z)=\frac1{1+z^2}\bigl(x(1+iz),y(1+iz)\bigr),
 \qquad x^2+y^2+z^2=1,
 \label{eq:whitney-sphere-dimension-two}
\end{equation}
其中 $x,y,z$ 是实数。

$0$ 处切空间的两个分支为
\[
\begin{aligned}
 V_+&=\operatorname{span}_{\mathbb R}
 \left\{(1+i)\frac\partial{\partial z_1},
        (1+i)\frac\partial{\partial z_2}\right\},&&z=1,\\
 V_-&=\operatorname{span}_{\mathbb R}
 \left\{(1-i)\frac\partial{\partial z_1},
        (1-i)\frac\partial{\partial z_2}\right\},&&z=-1.
\end{aligned}
\]
容易看出 $V_++V_-=\mathbb C^2$。以 $e_1,e_2$ 表示由如下次序给出的配备：
\[
 e_1=\{V_+,V_-\},\qquad e_2=\{V_-,V_+\},
\]
并以
\[
 L_\epsilon^{e_1},\qquad L_\epsilon^{e_2},\qquad\epsilon>0
\]
表示相应的拉格朗日手术。二者的拓扑类型都是环面 $T^2$。事实表明，对适当的
$\epsilon_1$，$L_\epsilon^{e_1}$哈密顿同痕于标准乘积环面
$S^1\times S^1\subset\mathbb C^2$，而 $L_\epsilon^{e_2}$ 并非如此。后一个环面已被证明同构于\emph{切卡诺夫环面} (Che96b)、(EP97)；其定义稍后解释。
\stfhIndex{I075}{qie^ka^nuo^fu^huan^mian}{切卡诺夫环面（Chekanov torus）}
\end{example}

\subsection{拉格朗日悬垂}
\stfhIndex{I191}{la^ge^lang^ri^xuan^chui}{拉格朗日悬垂（Lagrangian suspension）}
\label{subsec:lagrangian-suspension}

本小节以 $(X,\omega)$ 表示一般辛流形，而不用字母 $M$；还以 $Y$ 表示闭流形。给定 $Y$ 到 $(X,\omega)$ 中的拉格朗日嵌入时，我们往往不区分 $Y$ 与这个嵌入在 $X$ 中的像。

设 $\psi:[0,1]\times Y\to X$ 是一族嵌入 $\psi_t:Y\to X$ 的同痕。由
\cref{ex:smooth-primitive-exact-isotopy}，可以找到函数
$h:[0,1]\times Y\to\mathbb R$，使
$\psi^*\omega=dt\wedge dh$，即
\[
 \psi^*\omega+dt\wedge(-dh)=0.
\]
由\cref{thm:exact-lagrangian-isotopy-extension}，可以找到哈密顿量
$H:[0,1]\times X\to\mathbb R$，使
\begin{equation}
 \psi(t,x)=\phi_H^t(\psi(0,x))=\phi_H^t(\psi_0(x)).
 \label{eq:suspension-hamiltonian-isotopy}
\end{equation}
（给定正合拉格朗日同痕 $(h,\psi)$ 后，总可以在同痕的支集之外调整 $H$，使 $H$ 满足\emph{归一化}条件
\[
 \int_XH_t\,\omega^n=0.\text{）}
\]
此外，若同痕边界平坦，则 $h,H$ 也可取得边界平坦。也就是说，可以假设存在 $\epsilon_0>0$ 使
\begin{equation}
 H\equiv0\qquad\text{对 }t\in[0,\epsilon_0]\cup[1-\epsilon_0,1].
 \label{eq:boundary-flat-hamiltonian}
\end{equation}

若正合拉格朗日同痕 $\psi:[0,1]\times Y\to X$ 的相伴配对 $(h,\psi)$ 满足
\[
 \psi(0,x)=\psi(1,x),\qquad h(0,x)=h(1,x)
\]
对所有 $x\in Y$ 成立，则称它为\emph{正合拉格朗日环路}。

\begin{definition}\label{def:lagrangian-suspension}
设 $\psi:[0,1]\times Y\to X$ 是正合拉格朗日环路。与 $(h,\psi)$ 相伴的\emph{拉格朗日悬垂}是拉格朗日嵌入
\stfhIndex{I191}{la^ge^lang^ri^xuan^chui}{拉格朗日悬垂（Lagrangian suspension）}
\[
 \iota_{(h,\psi)}:S^1\times Y\longrightarrow
 \mathbb R\times S^1\times X=T^*S^1\times X,
\]
其定义为
\begin{equation}
 \iota_{(h,\psi)}(\theta,y)
 =(\theta,-h(\theta,x),\psi(\theta,x)).
 \label{eq:lagrangian-suspension}
\end{equation}
\end{definition}

下面陈述一个命题，它是拉格朗日悬垂构造及其应用中的关键成分。

考虑哈密顿路径的同痕，即二参数\termHamiltonianDiffeomorphism{}族
\[
 (s,\theta)\in[0,1]\times S^1\longmapsto
 \psi_s^\theta\in\operatorname{Ham}(X,\omega).
\]
以 $H=H(s,\theta,x)$ 表示生成 $\theta$-同痕
$\theta\mapsto\psi_s^\theta$ 的归一化哈密顿量，以
$K=K(s,\theta,x)$ 表示生成 $s$-同痕 $s\mapsto\psi_s^\theta$ 的哈密顿量。换言之，
\[
 \frac{\partial\psi_s^\theta}{\partial\theta}\circ(\psi_s^\theta)^{-1}=X_H,
 \qquad
 \frac{\partial\psi_s^\theta}{\partial s}\circ(\psi_s^\theta)^{-1}=X_K.
\]
回顾（2.5.34）中对归一化哈密顿量 $H$ 成立的恒等式
\begin{equation}
 \frac{\partial H}{\partial s}
 =\frac{\partial K}{\partial\theta}-\{K,H\}
 =\frac{\partial K}{\partial\theta}+\{H,K\}.
 \label{eq:two-parameter-hamiltonian-identity}
\end{equation}

\begin{proposition}\label{prop:suspension-hamiltonian-isotopic-product}
设 $(h,\psi)$ 是正合拉格朗日环路，并假设 $\psi$ 的像包含于
$U\subset X$。在 $T^*S^1\times X$ 上配备辛形式
\[
 d\theta\wedge db+\omega_X,
\]
其中 $(\theta,b)$ 是 $T^*S^1$ 的典范坐标。则拉格朗日嵌入
$\iota_{(h,\psi)}:S^1\times Y\to T^*S^1\times X$ 可由支集包含于
$T^*S^1\times U$ 的哈密顿流同痕到嵌入
$o_{S^1}\times i_Y:S^1\times Y\subset T^*S^1\times X$。
\end{proposition}

采用 (Oh07) 附录中的证明，来证明如下定理。

\begin{theorem}\label{thm:two-parameter-suspension-exact}
设 $Y\subset X$ 是拉格朗日嵌入，并考虑如上的二参数\termHamiltonianDiffeomorphism{}族
$\psi_s^t$。由
\begin{equation}
 f(s,\theta,x)
 =(\theta,-H(s,\theta,\psi_s^\theta(x)),\psi_s^\theta(x))
 \label{eq:two-parameter-suspension-map}
\end{equation}
定义的同痕
$f:[0,1]\times S^1\times Y\to T^*S^1\times X=:M$，在
\cref{def:exact-lagrangian-isotopy}的意义下是正合拉格朗日同痕。此外，
\begin{equation}
 i_s^*\alpha=d(K\circ f)_s.
 \label{eq:suspension-exact-primitive}
\end{equation}
\end{theorem}

需要说明，此处 $s$ 取代\cref{def:exact-lagrangian-isotopy}中 $t$ 的角色，而
$[a,b]\times Y$ 取代其中 $Y$ 的角色。

\begin{proof}
先计算
\[
 f^*(d\theta\wedge db+\omega_X).
\]
因为 $o_{S^1}\times i_Y\subset T^*S^1\times X$ 是拉格朗日嵌入，而
$\psi_s^\theta$ 是辛的，容易验证对 $[0,1]\times(S^1\times Y)$ 上某个 $1$-形式 $\alpha$ 有
\[
 f^*(d\theta\wedge db+\omega_X)=\alpha\wedge ds.
\]
并且
\[
 \alpha=\frac\partial{\partial s}\mathbin{\lrcorner}
 f^*(d\theta\wedge db+\omega_X).
\]
下面计算它。对一般切向量 $\xi\in T(S^1\times Y)$，先有
\[
 \alpha(\xi)=(d\theta\wedge db+\omega_X)
 \left(\frac{\partial f}{\partial s},df(\xi)\right).
\]
直接计算得到
\[
\begin{aligned}
 df\left(\frac\partial{\partial s}\right)
 &=\left(-\frac{\partial H}{\partial s}-\{H,K\}\right)
      \circ\psi_s^\theta\frac\partial{\partial b}
      \oplus X_K\circ\psi_s^\theta\\
 &=-\frac{\partial K}{\partial\theta}\circ\psi_s^\theta
      \frac\partial{\partial b}\oplus X_K\circ\psi_s^\theta,\\
 df\left(\frac\partial{\partial\theta}\right)
 &=\frac\partial{\partial\theta}\oplus
   -\frac{\partial H}{\partial\theta}\circ\psi_s^\theta
       \frac\partial{\partial b}\oplus X_H\circ\psi_s^\theta,\\
 df(\eta)
 &=-d_XH(d\psi_s^\theta(\eta))\frac\partial{\partial b}
      \oplus d\psi_s^\theta(\eta),
\end{aligned}
\]
其中 $\eta$ 是投影到 $[0,1]\times S^1$ 后为零的向量场。第一个等式使用恒等式
$dH(X_K)=\{H,K\}$，第二个等式使用（2.5.34）。由这些公式，得到
\[
\begin{aligned}
 i_s^*\alpha
 &=i_s^*\left(\frac\partial{\partial s}\mathbin{\lrcorner}
      (d\theta\wedge db+\omega_X)\right)\\
 &=\left(\left(\frac{\partial K}{\partial\theta}\circ f_s\right)
      d\theta\circ Tf_s+d_XK\circ Tf_s\right)\\
 &=f_s^*\left(\frac{\partial K}{\partial\theta}d\theta+d_XK\right)
 =f_s^*dK=d(K\circ f)_s.
\end{aligned}
\]
所以 $i_s^*\alpha$ 正合，$f$ 因而在\cref{def:exact-lagrangian-isotopy}的意义下定义了正合拉格朗日同痕；同时也证明了\eqref{eq:suspension-exact-primitive}。这就完成了证明。
\end{proof}

下面给出一种从任意给定正合拉格朗日同痕构造正合拉格朗日环路的简单方法。先引入同痕的\emph{时间反演}。

\begin{definition}\label{def:isotopy-time-reversal}
设 $Y\subset X$ 是拉格朗日子流形，$(h,\psi)$ 是正合拉格朗日同痕。$(h,\psi)$ 的\emph{时间反演}是配对 $(\widetilde h,\widetilde\psi)$：
\[
 \widetilde\psi(t,x)=\psi(1-t,x),
 \qquad
 \widetilde h(t,x)=-h(1-t,x).
\]
\end{definition}

显然，对边界平坦的 $(h,\psi)$，拼接同痕
\[
 (h,\psi)*(\widetilde h,\widetilde\psi)
\]
定义\emph{光滑}闭嵌入
$S^1(2)\times Y\to\mathbb R\times S^1(2)\times X$；这里把 $S^1$ 表成
$\mathbb R/2\mathbb Z$。称这个拼接同痕为同痕 $(h,\psi)$ 的\emph{奇双重}，并记作
\stfhIndex{I218}{qi^shuang^chong}{奇双重（odd-double）}
$(h^{\mathrm{od}},\psi^{\mathrm{od}})$。

\begin{definition}\label{def:odd-double-suspension}
正合拉格朗日同痕
$(h,\psi):[0,1]\times Y\to\mathbb R\times X$ 的\emph{奇双重悬垂}，是嵌入
\[
 \iota_{(h,\psi)}:S^1(2)\times Y\longrightarrow T^*S^1(2)\times X,
\]
其中
\begin{equation}
 \iota_{(h,\psi)}(\theta,x)=
 \begin{cases}
  (\theta,-h(\theta,x),\psi(\theta,x)),&\theta\in[0,1],\\
  (\theta,h(2-\theta,x),\psi(2-\theta,x)),&\theta\in[1,2].
 \end{cases}
 \label{eq:odd-double-suspension}
\end{equation}
\end{definition}

还用 $h^{\mathrm{od}}$ 表示哈密顿量 $h$ 的显然双重 $h*\widetilde h$。

这一奇双重悬垂构造，在 (Vi06) 和 (Oh07) 中证明拓扑哈密顿量的唯一性结果时起重要作用。

\subsection{切卡诺夫悬垂}
\label{subsec:chekanov-suspension}

本小节解释尤里·切卡诺夫（Yuri Chekanov）(Che96b) 引入的另一类悬垂；它只适用于
$\mathbb R^{2n}$ 中的拉格朗日子流形。

考虑浸没 $i_n:S^1\times\mathbb R^n\to\mathbb R^{n+1}$：
\[
 i_n(\theta,q_1,\ldots,q_n)
 =(e^{q_1}\cos(2\pi\theta),e^{q_1}\sin(2\pi\theta),q_2,\ldots,q_n),
\]
它定义了到 $(\mathbb R^2\setminus\{0\})\times\mathbb R^{n-1}$ 上的微分同胚。再以
\[
 i_n^*:T^*((\mathbb R^2\setminus\{0\})\times\mathbb R^{n-1})
 \longrightarrow T^*(S^1\times\mathbb R^n)
\]
表示 $i_n$ 的余微分。于是得到辛嵌入
\[
 I_n=(i_n^*)^{-1}:T^*(S^1\times\mathbb R^n)
 \longrightarrow T^*\mathbb R^{n+1}\cong\mathbb R^{2n+2}.
\]

现在设 $L\subset\mathbb R^{2n}\cong T^*\mathbb R^n$ 是拉格朗日子流形，$a$ 是任意实数。以 $N_a$ 表示平行 $1$-形式 $a\,d\theta$ 的图像，即在 $T^*S^1$ 的典范坐标中
\[
 \{(\theta,\alpha)\in T^*S^1\mid\alpha=a\}.
\]
把 $N_a\times L$ 看作
$T^*S^1\times T^*\mathbb R^n\cong T^*(S^1\times\mathbb R^n)$ 的拉格朗日子流形。于是 $L$ 的\emph{切卡诺夫悬垂}定义为
\[
 \Theta_a(L)=I_n(N_a\times L)
 \subset T^*(\mathbb R^{n+1})\cong\mathbb R^{2n+2}.
\]
切卡诺夫 (Che96b) 利用这个悬垂得到如下有趣结果。

\begin{theorem}[切卡诺夫]\label{thm:chekanov-suspension-torus}
设 $S^1\subset\mathbb R^2$ 是半径为 $1$ 的标准圆，并考虑悬垂
$\Theta_0(S^1)\subset\mathbb R^4$。则 $\Theta_0(S^1)$ 是 $\mathbb R^4$ 中的拉格朗日环面，但它不与标准环面
$S^1\times S^1$哈密顿同痕。
\end{theorem}

切卡诺夫 (Che96b) 最初通过研究埃克兰--霍费尔容量的一个变体证明了这一结果。随后埃利亚什伯格与波尔捷罗维奇 (EP97) 使用单点开格罗莫夫--威腾不变量的一个版本区分了这两个拉格朗日环面。

\section{典范关系与自然边界条件}
\label{sec:canonical-relation-natural-boundary}

本节研究一般辛流形上的最小作用量原理，并在开弦和闭弦两种情境中，为多值作用量泛函的变分理论建立适当框架。

正如第 1 章中余切丛上的经典力学那样，考虑光滑道路集合
\[
 \mathcal PM=\{\gamma:[0,1]\to M\}.
\]
有自然的评价映射
$\operatorname{ev}_0,\operatorname{ev}_1:\mathcal PM\to M$，由
$\operatorname{ev}_i(\gamma)=\gamma(i)$（$i=0,1$）给出。把它们合并，定义映射
\[
 \operatorname{ev}=(\operatorname{ev}_0,\operatorname{ev}_1):
 \mathcal PM\longrightarrow M\times M.
\]

下面的典范关系概念在拉尔斯·赫尔曼德（Lars Hörmander）的傅里叶积分算子演算中曾起重要作用；参见 (Hor71) 和 (GS77)。

\begin{definition}[典范关系]\label{def:canonical-relation}
\stfhIndex{I073}{dian^fan^guan^xi}{典范关系（canonical relation）}
若子集 $C\subset M\times M$ 关于 $M\times M$ 上的辛形式
$-\omega\oplus\omega$ 是拉格朗日子流形，则称 $C$ 为\emph{典范关系}。
\end{definition}

下面将说明，典范关系在变分理论中扮演自然边界条件的角色。

在 $\mathcal PM$ 上定义\emph{作用 $1$-形式} $\alpha$：
\stfhIndex{I053}{zuo^yong^1^xing^shi}{作用 1-形式（action one-form）}
\begin{equation}
 \alpha(\gamma)(\xi)=\int_0^1\omega(\dot\gamma(t),\xi(t))\,dt.
 \label{eq:path-space-action-one-form}
\end{equation}
下面的公式揭示典范关系与最小作用量原理之间的联系。

\begin{proposition}\label{prop:action-one-form-differential}
设 $\alpha$ 是 $\mathcal PM$ 上的作用 $1$-形式，则
\begin{equation}
 d\alpha=\operatorname{ev}_1^*\omega-\operatorname{ev}_0^*\omega.
 \label{eq:action-one-form-boundary-differential}
\end{equation}
（这里把 $\mathcal PM$ 看作无限维流形。）
\end{proposition}

\begin{proof}
须证明对任意 $\gamma\in\mathcal PM$ 及
$\xi,\eta\in T_\gamma(\mathcal PM)$，
\begin{equation}
 d\alpha(\gamma)(\xi,\eta)
 =\omega(\xi(1),\eta(1))-\omega(\xi(0),\eta(0)).
 \label{eq:action-one-form-boundary-values}
\end{equation}
为计算 $d\alpha(\gamma)(\xi,\eta)$，选取二参数道路族
\[
 V:(-\epsilon,\epsilon)\times(-\epsilon,\epsilon)\longrightarrow\mathcal PM,
\]
满足
\[
 V(0,0)=\gamma,\qquad
 \left.\frac{\partial V}{\partial u}\right|_{(0,0)}=\xi,\qquad
 \left.\frac{\partial V}{\partial s}\right|_{(0,0)}=\eta.
\]
也以 $V=V(u,s,t)$ 表示相应映射
\[
 V:(-\epsilon,\epsilon)\times(-\epsilon,\epsilon)\times[0,1]
 \longrightarrow M.
\]
于是
\begin{equation}
 d\alpha(\gamma)(\xi,\eta)
 =\left.\left[
 \frac\partial{\partial u}\alpha\left(\frac{\partial V}{\partial s}\right)
 -\frac\partial{\partial s}\alpha\left(\frac{\partial V}{\partial u}\right)
 \right]\right|_{(0,0)}.
 \label{eq:path-space-exterior-derivative-definition}
\end{equation}
这里使用了恒等式
\[
 \left[\frac{\partial V}{\partial u},\frac{\partial V}{\partial s}\right]
 =V_*\left[\frac\partial{\partial s},\frac\partial{\partial u}\right]=0.
\]

\begin{remark}\label{rem:path-space-differential-forms}
严格说来，在道路空间 $\mathcal PM$ 上讨论微分形式及其外微分时需要更加谨慎。这里建议读者把公式\eqref{eq:path-space-exterior-derivative-definition}当作
$\mathcal PM$ 上 $2$-形式 $d\alpha$ 的定义。关于道路空间上微分形式的更系统讨论，参见 (GJP91) 和 (Chen73)。
\end{remark}

由作用 $1$-形式\eqref{eq:path-space-action-one-form}的定义，
\[
 \frac\partial{\partial u}\alpha\left(\frac{\partial V}{\partial s}\right)
 =\int_0^1\frac\partial{\partial u}
 \left(\omega\left(\frac{\partial V}{\partial t},
                         \frac{\partial V}{\partial s}\right)\right)dt.
\]
现固定任意\emph{无挠辛联络} $\nabla$，则
\[
\begin{aligned}
 \frac\partial{\partial u}
 \left(\omega\left(\frac{\partial V}{\partial t},
                         \frac{\partial V}{\partial s}\right)\right)
 &=\omega\left(\nabla_u\frac{\partial V}{\partial t},
                    \frac{\partial V}{\partial s}\right)\\
 &\quad+\omega\left(\frac{\partial V}{\partial t},
                    \nabla_u\frac{\partial V}{\partial s}\right).
\end{aligned}
\]
由 $\nabla$ 的无挠条件，
\[
 \nabla_u\frac{\partial V}{\partial s}
 =\nabla_s\frac{\partial V}{\partial u}.
\]
所以第二项会与\eqref{eq:path-space-exterior-derivative-definition}中求
$(\partial\alpha/\partial s)(\partial V/\partial u)$ 时出现的类似项抵消。于是只需计算
\[
\begin{aligned}
 \omega\left(\nabla_u\frac{\partial V}{\partial t},
                  \frac{\partial V}{\partial s}\right)
 &=\omega\left(\nabla_t\frac{\partial V}{\partial u},
                  \frac{\partial V}{\partial s}\right)\\
 &=\frac\partial{\partial t}
    \left(\omega\left(\frac{\partial V}{\partial u},
                          \frac{\partial V}{\partial s}\right)\right)
 -\omega\left(\frac{\partial V}{\partial u},
                  \nabla_t\frac{\partial V}{\partial s}\right).
\end{aligned}
\]
交换 $u,s$ 的角色并使用 $\nabla$ 的辛性质，得到
\[
\begin{aligned}
 &\frac\partial{\partial u}
   \left(\omega\left(\frac{\partial V}{\partial t},
                        \frac{\partial V}{\partial s}\right)\right)
 -\frac\partial{\partial s}
   \left(\omega\left(\frac{\partial V}{\partial t},
                        \frac{\partial V}{\partial u}\right)\right)\\
 &=2\frac\partial{\partial t}
   \left(\omega\left(\frac{\partial V}{\partial u},
                        \frac{\partial V}{\partial s}\right)\right)
 +\omega\left(\frac{\partial V}{\partial s},
                 \nabla_t\frac{\partial V}{\partial u}\right)
 -\omega\left(\frac{\partial V}{\partial u},
                 \nabla_t\frac{\partial V}{\partial s}\right)\\
 &=2\frac\partial{\partial t}
   \left(\omega\left(\frac{\partial V}{\partial u},
                        \frac{\partial V}{\partial s}\right)\right)
 -\frac\partial{\partial t}
   \left(\omega\left(\frac{\partial V}{\partial u},
                        \frac{\partial V}{\partial s}\right)\right)\\
 &=\frac\partial{\partial t}
   \left(\omega\left(\frac{\partial V}{\partial u},
                        \frac{\partial V}{\partial s}\right)\right).
\end{aligned}
\]
因此
\[
\begin{aligned}
 d\alpha(\gamma)(\xi,\eta)
 &=\int_0^1\left.
   \frac\partial{\partial t}
   \left(\omega\left(\frac{\partial V}{\partial u},
                        \frac{\partial V}{\partial s}\right)\right)
   \right|_{(0,0)}dt\\
 &=\int_0^1\frac\partial{\partial t}\omega(\xi(t),\eta(t))\,dt\\
 &=\omega(\xi(1),\eta(1))-\omega(\xi(0),\eta(0)),
\end{aligned}
\]
这正是\eqref{eq:action-one-form-boundary-values}。
\end{proof}

对给定子集 $C\subset M\times M$，记
\[
 \mathcal P_CM=\operatorname{ev}^{-1}(C)
 =\{\gamma\in\mathcal PM\mid(\gamma(0),\gamma(1))\in C\}.
\]
仍用同一字母 $\alpha$ 表示它在 $\mathcal P_CM$ 上的显然限制。由公式\eqref{eq:action-one-form-boundary-differential}立即得到如下推论，它把典范关系刻画为自然边界条件。

\begin{corollary}\label{cor:canonical-relation-natural-boundary}
作用 $1$-形式 $\alpha$ 在 $\mathcal P_CM$ 上闭，当且仅当
$C\subset M\times M$ 是典范关系。
\end{corollary}

这里给出几个重要的典范关系例子；它们会出现在后续关于弗洛尔理论的各章中。

\begin{example}\label{ex:canonical-relations}
\begin{enumerate}[label=(\roman*)]
\item \textbf{自由环路。}对角线 $\Delta\subset M\times M$ 是典范关系。相应道路空间 $\mathcal P_\Delta M$ 可自然等同于自由环路空间
\[
 \mathcal LM=\{\gamma:S^1\to M\mid S^1=\mathbb R/\mathbb Z\}.
\]
$\alpha$ 在 $\mathcal P_\Delta M$ 上的零点集是常值道路的集合，因而就是
$\mathcal P_\Delta M$ 中的一份 $M\hookrightarrow\mathcal P_\Delta M$。

\item \textbf{扭曲环路。}辛微分同胚 $\phi:M\to M$ 的图像
\[
 \Delta_\phi:=\{(x,\phi(x))\mid x\in M\}
\]
是典范关系。相应道路空间是
\[
 \mathcal L_\phi M
 =\{\gamma\in\mathcal PM\mid\phi(\gamma(0))=\gamma(1)\},
\]
其元素称为\emph{扭曲环路}。事实上，可以把一个 $\phi$-扭曲环路等同于映射圆
\[
 M_\phi:=[0,1]\times M/{\sim}\longrightarrow S^1
\]
的一个截面，其中等价关系 $\sim$ 由
\[
 (0,x)\sim(1,x')\quad\Longleftrightarrow\quad x'=\phi(x)
\]
给出。这时 $\alpha$ 的零点集无非是 $\operatorname{Fix}\phi$ 中的常值环路。

\item \textbf{拉格朗日边界条件。}考虑形如
\[
 C=L_0\times L_1
\]
的典范关系，其中 $L_0,L_1$ 是 $M$ 的拉格朗日子流形。$\alpha$ 的零点集对应于交集 $L_0\cap L_1$。
\end{enumerate}
\end{example}

\section{生成函数与维泰博不变量}
\stfhIndexRangeStart{I136}{sheng^cheng^han^shu}{生成函数（generating function）}
\label{sec:generating-functions-viterbo-invariants}

\emph{生成函数}在哈密顿力学和微局部分析中起重要作用，尤其是在赫尔曼德的傅里叶积分算子演算 (Hor71) 中。生成函数在经典力学中的应用参见 (Go80)，它在傅里叶积分算子理论中的用法参见 (Hor71)、(GS77)。从几何观点看，给定流形 $N$ 上的一个函数族，在适当的横截性条件下，便可在 $T^*N$ 中得到一个自然的（浸入）拉格朗日子流形。这一构造可以纳入拉格朗日子流形演算中更一般的\emph{推前}运算，成为后者的特例。由于深谷范畴中态射与运算的函子性构造，这种拉格朗日子流形演算近来引起了广泛关注 (WW10)。本节用生成函数的构造来说明这种演算。

先给出推前拉格朗日子流形运算所依据的线性代数引理。回顾
\cref{lem:linear-symplectic-reduction}：对辛向量空间 $(V,\Omega)$ 的任意余迷向子空间 $C$，商空间 $C/C^\Omega$ 带有由 $\Omega$ 诱导的典范辛形式。（见\cref{lem:linear-symplectic-reduction}。）下面的引理说明，在什么条件下，拉格朗日子空间
$L\subset(V,\Omega)$ 会在 $C/C^\Omega$ 上典范地诱导一个拉格朗日子空间。

\begin{lemma}\label{lem:lagrangian-coisotropic-reduction}
设 $C\subset(V,\Omega)$ 是给定余迷向子空间。考虑拉格朗日子空间
$L\subset(V,\Omega)$ 以及交集 $L\cap C$。假设
\begin{equation}
 L+C=V.
 \label{eq:lagrangian-coisotropic-transversality}
\end{equation}
则 $\pi_C(L\cap C)\subset C/C^\Omega$ 是拉格朗日的。
\end{lemma}

\begin{proof}
$\pi_C(L\cap C)$ 在 $C/C^\Omega$ 中总是迷向的。再者，
\eqref{eq:lagrangian-coisotropic-transversality}蕴含
$L\cap C^\Omega=L^\Omega\cap C^\Omega=0$，所以
$\dim(L\cap C)=\dim\pi_C(L\cap C)$。还须证明
$\dim(L\cap C)=\frac12\dim(C/C^\Omega)$。因为
$\dim C+\dim C^\Omega=2n\;(=\dim V)$，有
\[
 \dim(C/C^\Omega)=2n-2\dim C^\Omega.
\]
另一方面，
\[
\begin{aligned}
 \dim(L\cap C)
 &=\dim L+\dim C-2n
 =\dim C-n=(2n-\dim C^\Omega)-n\\
 &=n-\dim C^\Omega
 =\frac12\dim(C/C^\Omega),
\end{aligned}
\]
故得证。
\end{proof}

现在考虑纤维丛 $\pi:E\to M$，以及纤维丛 $E$ 上的光滑函数
$S:E\to\mathbb R$。考虑拉格朗日子流形
$\operatorname{Graph}(dS)\subset T^*E$。我们希望把
$\operatorname{Graph}(dS)$ 推前到 $T^*N$ 中的拉格朗日子流形。考虑典范竖直切丛
$VTE\subset TE$，从而有子丛
\begin{equation}
 (VTE)^\perp\subset T^*E,
 \label{eq:vertical-annihilator-bundle}
\end{equation}
其中 $(VTE)^\perp$ 是 $T^*E$ 的子丛，它在 $e$ 处的纤维
$(VT_eE)^\perp\subset T_e^*E$ 是 $VT_eE$ 的零化子。注意
$\dim(VT_eE)^\perp=\dim T_{\pi(e)}^*N$。选取
$VT_eE\subset T_eE$ 的一个补空间，记作 $HT_eE$，便有同构
\[
 \pi_h^*:T_{\pi(e)}^*N\longrightarrow(HT_eE)^*
 \longrightarrow(VTE)^\perp,
\]
其中第一个映射是水平投影
$\pi_h:HT_eE\to T_{\pi(e)}N$ 的对偶映射。

\begin{exercise}\label{ex:vertical-annihilator-isomorphism}
定义这个同构，并证明其定义与分裂
$T_eE=VT_eE\oplus HT_eE$ 中补空间 $HT_eE$ 的选取无关。
\end{exercise}

因此有自然映射
\begin{equation}
 \operatorname{Graph}(dS)\cap(VTE)^\perp\longrightarrow T^*N.
 \label{eq:generating-function-pushforward-map}
\end{equation}
以 $L_S\subset T^*N$ 表示它的像。为了更具体地描述 $L_S$，回顾 $S$ 的\emph{纤维导数} $d^vS:E\to(VTE)^*$，它定义为
\[
 d^vS(e)=dS(e)|_{VT_eE},\qquad e\in E.
\]
于是可改写
\[
 \operatorname{Graph}(dS)\cap(VTE)^\perp
 =\{(e,D^hS(e))\mid d^vS(e)=0\},
\]
其中
\[
 D^hS(e):=dS(e)|_{HT_eE}.
\]

\begin{proposition}\label{prop:generating-function-lagrangian-immersion}
假设截面 $d^vS:E\to(VTE)^*$ 横截于零截面 $o_{(VTE)^*}$，则：
\begin{enumerate}[label=(\arabic*)]
\item 纤维临界点集
\[
 \Sigma_S:=\{e\in E\mid d^vS(e)=0\}
\]
是光滑流形；
\item 浸入 $i_S:\Sigma_S\to T^*N$ 的像
\[
 L_S=\{(q,dS|_{\Sigma_S}(e))\mid e\in\Sigma_S,\ \pi(e)=q\}
\]
是拉格朗日的。
\end{enumerate}
若（浸入）拉格朗日子流形 $L\subset T^*N$ 可以如上由 $S$ 表示，则称 $S$ 为 $L$ 的\emph{生成函数}。
\stfhIndexRangeEnd{I136}{sheng^cheng^han^shu}{生成函数（generating function）}
\end{proposition}

\begin{proof}
由 $\Sigma_S$ 的定义，第一项来自横截性定理。

为证明第二项，把 $e$ 写成 $e=(q,\xi)$，其中 $\pi(e)=q$ 且
$\xi\in VT_eE$；则映射 $i_S:\Sigma_S\to T^*M$ 为
\[
 i_S(e)=\bigl(\pi_E(e),D^hS(e)\circ(\pi_h(e))^{-1}\bigr),
 \qquad e\in\Sigma_S.
\]
现在计算
\[
\begin{aligned}
 i_S^*(\theta)(e)(v)
 &=\theta(di_S(e)(v))\\
 &=D^hS(e)\circ(\pi_h(e))^{-1}(d\pi_E(v))\\
 &=D^hS(e)(v)=i_{\Sigma_S}^*dS(e)(v),
\end{aligned}
\]
其中 $\theta=p\,dq$ 是 $T^*M$ 上的刘维尔 $1$-形式。这证明
$i_S^*(\theta)=i_{\Sigma_S}^*dS$，从而 $i_S^*\omega_0=0$。还须证明 $i_S$ 是浸入。

为此，设 $e\in\Sigma_S$、$v\in T_e\Sigma_S$ 且 $di_S(e)(v)=0$。于是
$d\pi_E(e)(v)=0$，故 $v$ 切于 $\pi_E:E\to N$ 的纤维。另一方面，$v$ 也切于
$\Sigma_S$。沿 $v\in T_e\Sigma_S$ 对 $d^vS(e)=0$ 微分，得到
$d(d^vS)(e)(v)=0$。因为 $v$ 竖直，$di_S(e)(v)=0$ 蕴含
$d^v(D^hS)(e)(v)=0$，故 $d^v(d^vS)(e)(v)=0$。而
$d^vS:E\to(VTE)^*$ 横截于 $(T^vE)^*$ 的零截面，等价于对每个
$e\in\Sigma_S$，
\[
 d^v(d^vS)(e):VT_eE\longrightarrow(VT_eE)^*
\]
是同构。因此 $v=0$。这证明 $i_S$ 是浸入，命题得证。
\end{proof}

\begin{remark}\label{rem:generating-function-euclidean-case}
当 $N=\mathbb R^n$、$E=\mathbb R^n\times\mathbb R^m$，且
$S=S(q,\xi)$ 是 $\mathbb R^n\times\mathbb R^m$ 上的函数时，可以等同
\[
 d^vS=d_\xi S,\qquad D^hS=d_qS,
\]
并且
$\Sigma_S=\{(q,\xi)\mid d_\xi S(q,\xi)=0\}$。此外，
$VTE\cong\mathbb R^n\times\mathbb R^m\times\mathbb R^m$，而
$d^vS:E\to(VTE)^*$ 可写作
\[
 d^vS(q,\xi)=(q,\xi,d^vS(q,\xi)).
\]
所以，在 $e=(q,\xi)\in\Sigma_S$ 处，$d^vS$ 横截于零截面
\[
 VTE\cong\mathbb R^n\times\mathbb R^m\times\{0\}
 \subset\mathbb R^n\times\mathbb R^m\times(\mathbb R^m)^*
\]
等价于黑塞 $\operatorname{Hess}(S|_{E_q})$ 可逆。
\end{remark}

\cref{prop:generating-function-lagrangian-immersion}的证明立即给出如下推论。

\begin{corollary}\label{cor:generated-lagrangian-exact}
$T^*N$ 中由生成函数生成的任意拉格朗日浸入都是正合的。
\end{corollary}

事实表明，并非每个正合拉格朗日浸入都由生成函数生成。关于用生成函数表示正合拉格朗日子流形时所遇到的这一障碍，详见 (Lat91)。关于正合拉格朗日子流形的生成函数存在性，弗朗索瓦·洛当巴克（François Laudenbach）与让-克洛德·西科拉夫（Jean-Claude Sikorav）的如下定理——即 (LS85) 的定理 2.1——给出了一项重要结果。

\begin{theorem}[洛当巴克--西科拉夫]\label{thm:laudenbach-sikorav-gfqi}
设 $L\subset T^*N$ 是零截面 $o_N\subset T^*N$ 的哈密顿形变
$L=\phi_H^1(o_N)$。则存在有限维向量丛 $E\to N$ 以及光滑函数
$S:E\to\mathbb R$ 来表示 $\phi_H^1(o_N)$，而且 $S$ 在无穷远处是二次的。
\end{theorem}

这里，称 $S$ \emph{在无穷远处二次}，是指存在非退化的逐纤维二次型 $Q$，使得在 $E$ 的某个紧子集之外
\stfhSubIndex{I138}{sheng^cheng^han^shu}{生成函数（generating function）}{zai^wu^qiong^yuan^chu^er^ci}{在无穷远处二次（quadratic at infinity）}
$S(q,\xi)=Q(q,\xi)$。受这个定理启发，下文只讨论这一特殊类别的生成函数。依照维泰博 (Vi92)，把在无穷远处二次的生成函数简称为 GFQI。

给定 $E$ 上逐纤维二次型 $Q$，以 $\mathcal S_{(Q;E)}$ 表示所有满足如下条件的 GFQI 之集合：在 $E$ 的一个紧子集之外
$S(q,\xi)=Q(q,\xi)$。注意，当 $E\to N$ 是向量丛时，有典范同构
$VT_eE\cong E_{\pi(e)}$，从而可作等同
\[
 VTE\cong\pi_E^*E,
 \qquad
 (VTE)^*\cong\pi_E^*E^*,
\]
其中 $\pi_E:E\to N$ 是投影。定义
\[
 \mathcal S^0_{(Q;E)}
 :=\{S\in\mathcal S_{(Q;E)}\mid
 d^vS\text{ 横截于 }o_{\pi_E^*E^*}\}.
\]
容易看出，在紧开拓扑下，$\mathcal S^0_{(Q;E)}$ 是
$\mathcal S_{(Q;E)}$ 的稠密子集。

下面引入生成函数上的两种运算；它们不改变所生成的拉格朗日子流形。

\begin{definition}\label{def:generating-function-equivalences}
\begin{enumerate}[label=(\arabic*)]
\item \textbf{稳定化。}若
\stfhSubIndex{I139}{sheng^cheng^han^shu}{生成函数（generating function）}{wen^ding^hua}{稳定化（stabilization）}
$\widetilde S:E\oplus F\to\mathbb R$ 满足
\[
 \widetilde S(q,\xi,\eta)=S(q,\xi)+Q(q,\eta)+C,
\]
其中 $Q$ 是向量丛 $F\to N$ 上某个非退化逐纤维二次型，$C$ 是常数，则称
$\widetilde S$ 为 $S$ 的一个\emph{稳定化}。

\item \textbf{规范等价。}若存在保纤维微分同胚 $\Phi:E\to E$ 使
\stfhSubIndex{I137}{sheng^cheng^han^shu}{生成函数（generating function）}{gui^fan^deng^jia}{规范等价（gauge equivalence）}
\[
 S_2=S_1\circ\Phi,
\]
则称 $S_1,S_2:E\to\mathbb R$ \emph{规范等价}。
\end{enumerate}
\end{definition}

下面这个唯一性结果是 (Vi92) 和 (Th99) 中证明的一个基本定理；其证明请参见这些文献。

\begin{theorem}[维泰博]\label{thm:viterbo-generating-function-uniqueness}
\cref{thm:laudenbach-sikorav-gfqi}中的生成函数，在相差稳定化与规范等价的意义下唯一。
\end{theorem}

现在给定 GFQI $S:E\to\mathbb R$。把
$S\in\mathcal S_{(Q;E)}$ 的次水平集记作
\[
 E^\lambda:=\{e\in E\mid S(e)\le\lambda\}.
\]
当 $|\lambda|$ 充分大时，
\[
 E^\lambda=\{e\in E\mid S(e)\le\lambda\}
 =\{e\in E\mid Q(e)\le\lambda\}
\]
的同伦型与 $\lambda$ 无关。对这样的 $\lambda$，把 $E^\lambda$ 记作
$E^\infty$，把 $E^{-\lambda}$ 记作 $E^{-\infty}$。注意
\[
 (E^\infty,E^{-\infty})\simeq(D(E^-),S(E^-)),
\]
其中 $E^-$ 表示二次型 $Q$ 的负丛，$D(E^-)$ 与 $S(E^-)$ 分别表示与
$E^-$ 相伴的圆盘丛与球面丛。

托姆同构给出同构
\begin{equation}
 H^*(M)\cong H^*(D(E^-),S(E^-))
 \cong H^*(E^\infty,E^{-\infty}).
 \label{eq:gfqi-thom-isomorphism}
\end{equation}
在 (Vi92) 中，维泰博为每个
$(u,S)\in H^*(M)\times\mathcal S^0_{(Q;E)}$ 配置实数
\[
 c(u,S)=\inf\{\lambda\mid j_\lambda^*Tu\ne0\}.
\]
这里 $T$ 表示同构\eqref{eq:gfqi-thom-isomorphism}，而
\[
 j_\lambda^*:H^*(E^\infty,E^{-\infty})
 \longrightarrow H^*(E^\lambda,E^{-\infty})
\]
是由包含
\[
 j_\lambda:E^\lambda\longrightarrow E^\infty
\]
诱导的同态。

依照 (MiO95)，通过考虑任意闭子流形 $K\subset N$，我们推广这一构造；以
$\pi|_K:E|_K\to K$ 表示丛 $E$ 在 $K$ 上的限制。

\begin{definition}\label{def:viterbo-minimax-invariant}
给定 $S\in\mathcal S^0_{(Q;E)}$ 与 $K\subset N$。对每个
$0\ne u\in H^*(K)$，定义
\[
 c(S;u,K):=\inf\{\lambda\mid
 j_\lambda^*Tu\ne0\text{ 于 }H^*(E_K^\lambda,E_K^{-\infty})\text{ 中}\}.
\]
\end{definition}

下面的结果既可按维泰博在 (Vi92) 中的方法证明，也可像 (Oh97b) 那样用弗洛尔理论的方法证明。

\begin{theorem}[参见 (Vi92) 的定理 2.3]\label{thm:viterbo-invariant-properties}
设 $S\in\mathcal S^0_{(Q;E)}$ 是拉格朗日子流形
$L\subset T^*N$ 的生成函数，$K\subset N$ 是闭子流形，则有：
\begin{enumerate}[label=(\arabic*)]
\item $c(S;u,K)$ 是 $S|_{E_K}$ 的临界值。
\item 对每个 $0\ne u\in H^*(K)$，存在点
$x_u\in\nu^*K\cap L$，使
\[
 c(S;u,K)-c(S;v,K)=\int_\gamma\theta,
\]
其中 $\gamma$ 是 $L$ 中任意一条连接 $x_u$ 与 $x_v$ 的光滑路径。
\item 分别以 $1$ 与 $\mu_N$ 表示 $H^0(N)$ 和 $H^n(N)$ 的生成元。记
\[
 \gamma(L):=c(\mu_N,S)-c(1,S).
\]
则 $\gamma(L)\ge0$，并且 $\gamma(L)=0$ 当且仅当 $L=o_M$。
\item 在 $S$ 的紧开拓扑下，$S\mapsto c(S;u,K)$ 是连续函数。
\end{enumerate}
\end{theorem}

